\documentclass[11pt,a4paper]{article}

% --- packages -----------------------------------------------------------
\usepackage[utf8]{inputenc}
\usepackage[T1]{fontenc}
\usepackage[margin=1in]{geometry}
\usepackage{amsmath,amssymb,amsthm}
\usepackage{bm}
\usepackage{mathtools}        % for \boxed inside align*, \DeclarePairedDelimiter, etc.
\usepackage{booktabs}
\usepackage{xcolor}
\usepackage{hyperref}
\hypersetup{colorlinks=true, linkcolor=blue!50!black, citecolor=blue!50!black,
            urlcolor=blue!50!black}

% --- convenient macros --------------------------------------------------
\newcommand{\eps}{\varepsilon}
\newcommand{\epstilde}{\widetilde{\varepsilon}}
\newcommand{\bvec}[1]{\mathbf{#1}}
\newcommand{\E}{\bvec{E}}
\newcommand{\f}{\bvec{f}}
\newcommand{\h}{\bvec{h}}
\newcommand{\ri}{\bvec{r}}
\newcommand{\dthreer}{\,d^{3}r}
\newcommand{\chitwo}{\chi^{(2)}}
\newcommand{\chithree}{\chi^{(3)}}
\newcommand{\Dom}{\Delta\omega}

% spacing tweaks
\setlength{\parskip}{0.4em}

% =======================================================================
\begin{document}

\title{Cascaded $\chi^{(2)}{:}\chi^{(2)} \to \chi^{(3)}_\text{eff}$ \\
       in a Doubly-Resonant Dielectric Metasurface}
\author{}
\date{\today}
\maketitle

% =======================================================================
\section{Setting and goal}

The objective is to engineer a metasurface whose effective third-order
nonlinear susceptibility, \emph{driven only at the fundamental frequency}
$\omega_1$, exceeds the bulk $\chithree$ of the constituent material by
orders of magnitude.  The mechanism is the well-known \emph{cascade} of
two $\chitwo$ processes:
$\omega_1+\omega_1\to 2\omega_1$ followed by $2\omega_1-\omega_1\to\omega_1$.
When the second-harmonic (SH) channel is \emph{off-resonant} (or
phase-mismatched), the SH amplitude follows the fundamental adiabatically
and its back-action produces a Kerr-like response on the fundamental
\cite{Ostrovskii1967,DeSalvo1992,Stegeman1996}.

The metasurface platform we have in mind consists of an array of
high-index dielectric resonators (LiNbO$_3$, GaAs, GaP, Si) supporting two
co-localised resonances:
\begin{itemize}
  \item a fundamental resonance at $\omega_1$ -- a \emph{quasi}-bound
        state in the continuum (qBIC) with controlled radiative coupling
        $\gamma_1^{(\text{rad})}$ to free space, designed for
        in-coupling at $\omega_1$;
  \item a second resonance at $\omega_2 \approx 2\omega_1$ -- a
        symmetry-protected BIC (or strongly under-coupled qBIC) whose
        total decay $\gamma_2$ is dominated by intrinsic absorption,
        giving the largest accessible quality factor
        $Q_2 = \omega_2/\gamma_2$.
\end{itemize}

The relevant tensor of $\chitwo$ is patterned in space (via the geometry
of the meta-atoms, by periodic poling, or by symmetry-breaking
perturbations) so that the SHG overlap integral does not vanish
(Section~\ref{sec:tensor}).  High-$Q$ metasurface SHG of exactly this
kind has been demonstrated in AlGaAs and LiNbO$_3$
\cite{Carletti2018,Koshelev2020,Anthur2020}.  These experiments measure
the SH radiation leaving the sample and therefore probe the
$\chitwo$-driven up-conversion channel; the \emph{cascaded}
$\chithree$ regime is the complementary observable -- an
intensity-dependent phase shift on the FM in conditions where the SH
radiation is intentionally suppressed -- and has not, to our
knowledge, been isolated experimentally in a metasurface platform.
A literature audit before submission is warranted, since
$\chitwo$:$\chitwo$ cascading in waveguides and bulk PPLN
\cite{DeSalvo1992,Stegeman1996,Bache2007} is well-established and the
metasurface extension is a natural step.

% =======================================================================
\section{Conventions}\label{sec:conventions}

Several factors of $2$ and $3$ in the final answer come from book-keeping
in the field/polarisation conventions; we fix them once and for all.

\paragraph{Real fields.}
Physical real fields are written as a sum of complex envelopes,
\begin{equation}
  \mathcal{E}_i(\ri,t)=\tfrac12\sum_n\!\bigl[E_{n,i}(\ri,t)\,e^{-i\omega_n t}+\mathrm{c.c.}\bigr],\qquad
  \mathcal{P}_i(\ri,t)=\tfrac12\sum_n\!\bigl[P_{n,i}(\ri,t)\,e^{-i\omega_n t}+\mathrm{c.c.}\bigr],
\end{equation}
with $E_{n,i},P_{n,i}$ slowly varying compared with $1/\omega_n$.

\paragraph{Susceptibility tensors (Boyd convention).}
With this field convention,
\begin{equation}
  P_i^{(2)}(\omega_\sigma)=\eps_0\!\!\sum_{(p,q)}\!\chitwo_{ijk}(\omega_\sigma;\omega_p,\omega_q)\,E_j(\omega_p)\,E_k(\omega_q),
\end{equation}
where the sum runs over distinct ordered pairs of frequencies summing to
$\omega_\sigma$.  Indices are summed by Einstein convention.

For the two processes at hand,
\begin{align}
  \boxed{\;P_i^{(2)}(2\omega_1)=\eps_0\,\chitwo_{ijk}(2\omega_1;\omega_1,\omega_1)\,E_j(\omega_1)\,E_k(\omega_1)\;}
  &\qquad\text{(SHG)} \label{eq:PSHG}\\[4pt]
  \boxed{\;P_i^{(2)}(\omega_1)=2\eps_0\,\chitwo_{ijk}(\omega_1;2\omega_1,-\omega_1)\,E_j(2\omega_1)\,E_k^*(\omega_1)\;}
  &\qquad\text{(DFG)}\label{eq:PDFG}
\end{align}
The factor $2$ in DFG counts the orderings $(2\omega_1,-\omega_1)$ and
$(-\omega_1,2\omega_1)$.  \textbf{This is the factor missing in the
original note}, and dropping it violates Manley--Rowe (Section~\ref{sec:tcmt}).

The native Kerr response on the fundamental, for comparison, is
\begin{equation}
  P_i^{(3)}(\omega_1)=3\eps_0\,\chithree_{ijkl}(\omega_1;\omega_1,\omega_1,-\omega_1)\,
                     E_j(\omega_1)E_k(\omega_1)E_l^*(\omega_1),
\end{equation}
with the factor $3$ counting the distinct orderings of
$(\omega_1,\omega_1,-\omega_1)$.

We assume Kleinman symmetry (full permutation symmetry across all index
pairs), valid for both LiNbO$_3$ and GaAs at telecom-band fundamentals
well below the electronic gap.  In particular,
\begin{equation}
  \chitwo_{ijk}(2\omega_1;\omega_1,\omega_1)=\chitwo_{kij}(\omega_1;2\omega_1,-\omega_1)\equiv\chitwo_{ijk}(\ri),
\end{equation}
i.e.\ a single symmetric tensor $\chitwo_{ijk}(\ri)$ in what follows.

% =======================================================================
\section{Modes and energy normalisation}

Let $\f_n(\ri),\h_n(\ri)$ be the eigenmodes of the linear, lossless,
dispersive Maxwell problem,
\begin{equation}
  \nabla\!\times\!\nabla\!\times\!\f_n(\ri)
  =\frac{\omega_n^2}{c^2}\,\eps_r(\ri,\omega_n)\,\f_n(\ri),
\end{equation}
with $\f_n$ a complex vector field encoding the polarisation pattern of
mode $n$ (and the magnetic mode
$\h_n=\nabla\!\times\!\f_n/(i\omega_n\mu_0)=-i\nabla\!\times\!\f_n/(\omega_n\mu_0)$,
following Faraday's law $\nabla\!\times\!\E=-\partial_t\mathbf B$ under the
$e^{-i\omega t}$ convention of \S\ref{sec:conventions}).
Note that both the \emph{vectorial} nature of $\f_n$ and the dispersion
of $\eps_r$ are essential; the original note's scalar, dispersionless
$f(\ri)$ drops information needed for the tensorial overlap below.

The mode is energy-normalised so that
\begin{equation}
  \boxed{\;\frac12\!\int\!\!\bigl[\,\eps_0\,\epstilde_n(\ri)\,|\f_n|^2+\mu_0|\h_n|^2\,\bigr]\dthreer=1,
  \qquad
  \epstilde_n(\ri)\equiv\left.\frac{\partial(\omega\,\eps_r(\ri,\omega))}{\partial\omega}\right|_{\omega=\omega_n}\;}
\label{eq:Brillouin}
\end{equation}
(Brillouin's expression for the energy density of a time-harmonic field
in a dispersive medium \cite[\S 80]{LandauLifshitz}).  Defining the
modal amplitude $a_n(t)$ by
\begin{equation}
  E_{n,i}(\ri,t)=a_n(t)\,f_{n,i}(\ri),
\end{equation}
the time-averaged energy stored in mode $n$ is exactly $|a_n|^2$
(so $a_n$ has units of $\sqrt{\mathrm J}$ and $|a_n|^2$ has units of $\mathrm J$).

For non-dispersive lossless dielectrics one may replace
$\epstilde_n\to\eps_r(\ri,\omega_n)\equiv n^2(\ri,\omega_n)$.
\textbf{Do not use a single $n$ for both modes}: in LiNbO$_3$ at
$\lambda_1=1550$~nm, $n_o(\omega_1)\simeq 2.21$ and $n_o(2\omega_1)\simeq 2.26$
(a $\sim 2\%$ difference), and the magnitude of $\chitwo_{ijk}$ also
differs at the two frequencies via Miller's rule.

% =======================================================================
\section{Polarisation-tensor derivation: $\chithree_{\text{casc}}$ from $\chitwo\!\cdot G\!\cdot\chitwo$}\label{sec:tensorder}

Before introducing modal amplitudes it is instructive to derive
$\chithree_{\text{casc}}$ directly from the polarisation expansion and
Maxwell's equations, without invoking coupled-mode theory.  This route
makes the structural form
$\chithree_{\text{casc}}\sim\chitwo\!\cdot G(2\omega_1)\!\cdot\chitwo$
manifest, exposes a non-local kernel that single-mode TCMT silently
collapses, and shows that the temporal coupled-mode result of
Section~\ref{sec:tcmt} is the same calculation re-expressed in a
single-resonance approximation.

\paragraph{The SH field as the response to a $\chitwo$ source.}
At leading order in the nonlinearity the SH field is determined by
Maxwell's equations with the SHG polarisation \eqref{eq:PSHG} as a
source,
\begin{equation}
  \bigl[\nabla\!\times\!\nabla\!\times\,-\,k_2^2\,\eps_r(\ri,2\omega_1)\bigr]
  \mathbf E(2\omega_1)=k_2^2\,\mu_0 c^2\,\mathbf P^{\text{NL}}(2\omega_1),
  \qquad k_2\equiv 2\omega_1/c.
\label{eq:waveeq-SH}
\end{equation}
Define $\overline{\overline G}(\ri,\ri';2\omega_1)$ as the linear,
retarded, dispersive dyadic Green's function for the operator on the
left -- i.e.\ the field at $\ri$ produced by a unit polarisation source
at $\ri'$ in the same structure.  Substituting \eqref{eq:PSHG},
\begin{equation}
  \boxed{\;
    E_j(\ri,2\omega_1)
    = k_2^2\!\int\!d^3r'\;
      G_{jl}(\ri,\ri';2\omega_1)\,\chitwo_{lmn}(\ri')\,
      E_m(\ri',\omega_1)\,E_n(\ri',\omega_1).
  \;}
  \label{eq:Esh-Green}
\end{equation}
This is the \emph{only} step that goes beyond the susceptibility
tensors: it tells us how the SH field is geometrically related to its
$\chitwo$ source in the chosen structure (bulk medium, waveguide,
metasurface, photonic crystal, $\ldots$).

\paragraph{Cascaded polarisation at $\omega_1$.}
Substituting \eqref{eq:Esh-Green} into the DFG term \eqref{eq:PDFG}
gives, with summation over repeated indices,
\begin{equation}
  \boxed{\;
  P_i^{\text{casc}}(\ri,\omega_1)
  = 2\eps_0\,k_2^2\!\int\!d^3r'\;
    \underbrace{\chitwo_{ijk}(\ri)\,
                G_{jl}(\ri,\ri';2\omega_1)\,
                \chitwo_{lmn}(\ri')}_{\equiv\,\Xi^{\text{casc}}_{ikmn}(\ri,\ri')}
    E_m(\ri',\omega_1)\,E_n(\ri',\omega_1)\,E_k^*(\ri,\omega_1).
  \;}\label{eq:Pcasc-nonlocal}
\end{equation}
Two structural features of \eqref{eq:Pcasc-nonlocal} are worth
emphasising before specialising.

\textbf{The cascade is non-local.}  In general
$P_i^{\text{casc}}(\ri,\omega_1)$ depends on the FM field at a different
point $\ri'$ through the SH Green's function.  A genuinely \emph{local}
effective $\chithree(\ri)$ exists only when $G(\ri,\ri';2\omega_1)$ is
sufficiently sharply peaked at $\ri=\ri'$, or when the FM field varies
slowly on the support of $G$.  This is automatic in bulk media
(translational invariance shrinks the kernel to a
$\Delta k$-dependent local response) and in single-mode metasurface
unit cells, but \emph{not} in regimes where, e.g., the FM is a tightly
focused beam probing a weakly coupled lattice resonance.  In such
regimes the local-$\chithree$ description breaks down even though the
underlying cascaded mechanism is unchanged.

\textbf{The factor $2/3$ from the Boyd degeneracies is universal.}
Comparing \eqref{eq:Pcasc-nonlocal} with the native-Kerr definition
$P^{(3)}(\omega_1)=3\eps_0\chithree|E|^2 E$ shows that the cascaded
$\chithree$ inherits a factor $2/3$ relative to a naive product
$|\chitwo|^2/G$, simply because the DFG channel has degeneracy $2$
while the SPM channel has degeneracy $3$.  This shifts the answer in
\emph{both} the bulk and resonant limits below.

\paragraph{Two specialisations.}
Different physical contexts correspond to different forms of $G$:
\begin{itemize}\setlength{\itemsep}{3pt}
  \item \textbf{(a) Bulk, phase-mismatched.}  A translationally
    invariant medium has $G(\ri,\ri';2\omega_1)$ that depends only on
    $\ri-\ri'$; in $\mathbf k$-space it has a smooth dependence peaked
    at $|\mathbf k|=2\omega_1 n_2/c$.  For plane-wave fundamental
    fields with wave-vector mismatch
    $\Delta\mathbf k=\mathbf k_2-2\mathbf k_1\neq 0$ the integral in
    \eqref{eq:Pcasc-nonlocal} collapses to a local response.  The
    SVEA reduction \cite{DeSalvo1992,Stegeman1996,Boyd} proceeds as
    follows.  Writing the fields as
    $E(\omega_n,z)=A_n(z)\,e^{i k_n z}$ with slowly varying envelopes
    $A_n$, the wave equation \eqref{eq:waveeq-SH} at $2\omega_1$ with
    the SHG polarisation \eqref{eq:PSHG} as source reduces (after
    dropping $\partial_z^{2}A_2$ relative to $k_2\partial_z A_2$) to
    \cite[\S 2.4]{Boyd}
    \begin{equation}
      \frac{dA_2}{dz}\;=\;\frac{i\,\omega_2}{2\,n_2\,c}\,\chitwo\,A_1^{2}\,
                      e^{-i\Delta k z}.
      \label{eq:SVEA-A2}
    \end{equation}
    For undepleted pump ($A_1=$const) and $A_2(0)=0$ this integrates to
    \begin{equation}
      A_2(z)\;=\;\frac{\omega_2\,\chitwo\,A_1^{2}}{2\,n_2\,c\,\Delta k}\,
                \bigl(1-e^{-i\Delta k z}\bigr).
      \label{eq:A2-of-z}
    \end{equation}
    Substituting \eqref{eq:A2-of-z} into the DFG polarisation
    \eqref{eq:PDFG} and projecting the result back into the SVEA
    equation at $\omega_1$,
    $dA_1/dz=(i\omega_1/(2 n_1 c)) \cdot P^\text{NL}(\omega_1)/\eps_0
    \cdot e^{-i k_1 z}$, gives an intensity-dependent term
    $dA_1/dz\supset iK_\text{casc}|A_1|^{2}A_1$ with
    \begin{equation*}
      K_\text{casc}\;=\;\frac{\omega_1\,\omega_2}{2\,n_1\,n_2\,c^{2}\,\Delta k}
                       \,|\chitwo|^{2}\;\times\;(\text{phase factor}).
    \end{equation*}
    Identifying with the effective-Kerr SVEA equation
    $dA_1/dz=i\,(3\omega_1/(2 n_1 c))\,\chithree_\text{casc}\,|A_1|^{2}A_1$
    (where the factor of $3$ is the SPM Boyd degeneracy from \eqref{eq:PSHG}--\eqref{eq:PDFG})
    gives $\chithree_\text{casc}=K_\text{casc}\,2n_1c/(3\omega_1)$; the
    explicit $n_1$ from the FM projection cancels against the $n_1$ in
    $K_\text{casc}$, leaving only the SH refractive index:
    \begin{equation}
      \boxed{\;\chithree_{\text{casc,\,bulk}}(\omega_1)
              \;=\;-\,\frac{2\,\omega_1}{3\,c\,n_2\,\Delta k}\;
                   |\chitwo|^{2}\,\mathcal{F}(\Delta k\,L),\;}
      \label{eq:chi3casc-bulk}
    \end{equation}
    with $\mathcal{F}(\Delta k L)$ a dimensionless phase factor that
    averages to $\mathcal{F}\!\to\!1$ for $\Delta k L\!\gg\!1$
    (broadband bulk cascading; the regime exploited in soliton
    compression) and to $\mathcal{F}\!\to\!\Delta k L/2$ for
    $\Delta k L\!\ll\!1$ (where SHG conversion competes with the
    cascade).  Dimensionally $[\omega_1/(c\,\Delta k)]$ is dimensionless
    and $[|\chitwo|^{2}/n_2]=\mathrm{m}^{2}/\mathrm{V}^{2}$ as required;
    the explicit factor $2/3$ is the DFG-to-SPM Boyd degeneracy ratio
    inherited from \eqref{eq:Pcasc-nonlocal}.  Note that only $n_2$
    (the SH-mode index) survives in the end, paralleling the metasurface
    result \eqref{eq:chi3casc_explicit} below.  The
    cascaded contribution is sign-tunable through $\mathrm{sign}(\Delta k)$
    and bounded by the coherence length $\pi/\Delta k$ over which the
    cascade builds up coherently before reversing.  Equation
    \eqref{eq:chi3casc-bulk} is used in Section~\ref{sec:numerical} to
    estimate the bulk PPLN cascaded $\chithree$ for LiNbO$_3$.
  \item \textbf{(b) Single resonance (the metasurface case).}  Near an
    isolated SH quasi-normal mode at $\omega_2\approx 2\omega_1$ with
    vectorial profile $\f_2(\ri)$ and decay rate $\gamma_2$, the QNM
    expansion~\cite{Kristensen2020,Sauvan2013,Lalanne2018} gives, near
    the pole,
    \begin{equation}
      G_{jl}(\ri,\ri';2\omega_1)\;=\;
      \frac{c^2}{2\omega_2}\,
      \frac{f_{2,j}(\ri)\,f_{2,l}^*(\ri')}{\Dom+i\gamma_2/2}
      \;+\;(\text{non-resonant background}),
      \qquad\Dom=2\omega_1-\omega_2,
    \label{eq:G-QNM}
    \end{equation}
    with $\f_2$ in the Brillouin normalisation \eqref{eq:Brillouin}.
    The $c^2/(2\omega_2)$ prefactor is the standard residue of the
    one-pole QNM expansion in the wave-equation form
    $[\nabla\!\times\!\nabla\!\times\,-\,(\omega/c)^{2}\eps_r]^{-1}$
    (see~\cite[Eq.~(31)]{Sauvan2013}); we will not need its exact
    numerical value because the chain ending in
    \eqref{eq:chi3casc-tensor} is anchored against the TCMT route of
    Sections~\ref{sec:tcmt}--\ref{sec:adiab}, where the prefactor is
    fixed independently by the Brillouin energy normalisation.
    If the FM field is also dominated by a single mode,
    $E_m(\ri,\omega_1)=a_1\,f_{1,m}(\ri)$, substituting \eqref{eq:G-QNM}
    into the kernel
    $\Xi^{\text{casc}}_{ikmn}(\ri,\ri')=\chitwo_{ijk}(\ri)\,G_{jl}(\ri,\ri';2\omega_1)\,\chitwo_{lmn}(\ri')$
    of \eqref{eq:Pcasc-nonlocal} and using
    $E_m(\ri',\omega_1)E_n(\ri',\omega_1)=a_1^{2}f_{1,m}(\ri')f_{1,n}(\ri')$
    isolates two overlap integrals
    \begin{equation}
      M\equiv \!\int\!\chitwo_{ijk}(\ri)\,f_{2,i}^*\,f_{1,j}\,f_{1,k}\,\dthreer,
      \qquad
      \overline M\equiv \!\int\!\chitwo_{ijk}(\ri)\,f_{1,i}^*\,f_{2,j}\,f_{1,k}^*\,\dthreer,
      \label{eq:M-overline-M}
    \end{equation}
    related by Kleinman symmetry as $\overline M=M^*$ (the integrand
    of $\overline M$ is the complex conjugate of the integrand of $M$
    under the symmetry $\chitwo_{ijk}=\chitwo_{kij}$ of
    Section~\ref{sec:conventions}).  The integration over $\ri'$
    contracts $f_{2,l}^*(\ri')$ with $\chitwo_{lmn}(\ri')\,f_{1,m}(\ri')\,f_{1,n}(\ri')$
    -- precisely $M$ -- while the remaining factor at $\ri$,
    $\chitwo_{ijk}(\ri)\,f_{2,j}(\ri)\,E_k^*(\ri)=a_1^*\,\chitwo_{ijk}\,f_{2,j}\,f_{1,k}^*$,
    has the modal polarisation pattern of $\f_1(\ri)$ when projected
    appropriately.  The cascaded polarisation is therefore
    \begin{equation}
      P_i^{\text{casc}}(\ri,\omega_1)\;\propto\;
        \frac{M\,|a_1|^{2}\,a_1}{\Dom+i\gamma_2/2}\,
        \chitwo_{ijk}(\ri)\,f_{2,j}(\ri)\,f_{1,k}^{*}(\ri),
      \label{eq:Pcasc-modal}
    \end{equation}
    Kerr-like in the modal sense: it carries the SPM time-dependence
    $|a_1|^{2}a_1$ and inherits the single-pole Lorentzian
    $1/(\Dom+i\gamma_2/2)$ from $G$.  Projecting \eqref{eq:Pcasc-modal}
    onto $f_{1,i}^{*}(\ri)$ replaces the spatial factor by
    $\overline M=M^{*}$, so $\int f_{1,i}^{*}P_i^{\text{casc}}\dthreer
     \propto |M|^{2}\,|a_1|^{2}\,a_1/(\Dom+i\gamma_2/2)$; equating
    this to the modal projection of the native Kerr term
    $\int f_{1,i}^{*}P_i^{\text{nat}}\dthreer
     =3\eps_0\,\chithree_\text{casc}\,|a_1|^{2}\,a_1\!\int|\f_1|^{4}\dthreer$
    identifies the effective $\chithree$,
    \begin{equation}
      \boxed{\;\chithree_{\text{casc}}\;=\;
        -\,\frac{\eps_0\,\omega_1}{6}\cdot
        \frac{|M|^2}{(\Dom+i\gamma_2/2)\!\int\!|\f_1|^4\dthreer}\;}
      \label{eq:chi3casc-tensor}
    \end{equation}
    whose real part recovers \eqref{eq:chi3casc} exactly in the
    adiabatic limit $|\Dom|\gg\gamma_2$, and whose imaginary part is
    the cascaded two-photon-like loss derived in
    Section~\ref{sec:adiab}.  The overall prefactor $\eps_0\omega_1/6$
    depends on the QNM Green's-function normalisation in \eqref{eq:G-QNM}
    and on the Brillouin energy normalisation \eqref{eq:Brillouin};
    in the convention adopted here it is most cleanly fixed by the
    TCMT route of Sections~\ref{sec:tcmt}--\ref{sec:adiab}, where the
    same answer is obtained from the modal projection of Maxwell's
    equations directly.
\end{itemize}

\paragraph{Equivalence with the TCMT derivation.}
Using the resonant Green's function \eqref{eq:G-QNM} directly in
\eqref{eq:Esh-Green} gives $\mathbf E(\ri,2\omega_1)=a_2\f_2(\ri)$
with
\begin{equation}
  a_2\;\propto\;\frac{a_1^2}{\Dom+i\gamma_2/2}\,,
\end{equation}
which is precisely the steady state of the TCMT equation
\eqref{eq:a2}.  Matching prefactors recovers the explicit coupling
$\beta=(\eps_0\omega_1/4)\,M$ in \eqref{eq:beta}; the precise numerical
coefficient depends on the field/normalisation convention adopted in
Section~\ref{sec:conventions}.  \textbf{The temporal coupled-mode
theory is therefore the same calculation as \eqref{eq:Pcasc-nonlocal}
written in mode-amplitude language}, with $G(2\omega_1)$ replaced by
its single-pole Lorentzian.

\paragraph{What this reframing buys.}
\begin{itemize}\setlength{\itemsep}{3pt}
  \item The structural formula
    $\chithree_{\text{casc}}\sim\chitwo\!\cdot G(2\omega_1)\!\cdot\chitwo$
    makes clear that one can engineer the cascade through \emph{any}
    feature of $G$ at $2\omega_1$: phase mismatch, BIC, photonic-crystal
    band edge, polariton, slow light.  The high-$Q$ SH mode of
    Section~\ref{sec:tcmt} is one (very effective) choice, not the only
    one.
  \item The non-local kernel $\Xi^{\text{casc}}(\ri,\ri')$ becomes
    important whenever the FM field varies on the scale of the SH mode
    -- relevant for lattice resonances, surface-bound modes, or
    metasurfaces probed by tightly focused beams.  In those regimes
    the local single-mode picture used in
    Sections~\ref{sec:tcmt}--\ref{sec:tensor} should be replaced by a
    full evaluation of \eqref{eq:Pcasc-nonlocal}.
  \item The factor $2/3$ from Boyd degeneracies is the same in bulk
    and resonant cases -- a useful sanity check whenever comparing to
    a Stegeman-style bulk estimate~\cite{Stegeman1996}.
\end{itemize}

% =======================================================================
\section{Coupled-mode equations}\label{sec:tcmt}

Substituting the modal expansion into Maxwell's equations and projecting
onto $\f_n^*$ (slowly-varying-envelope approximation; e.g.\
\cite{SuhWangFan2004,Rodriguez2007}) gives, after adding phenomenological
loss rates and the input port for the FM:
\begin{align}
  \boxed{\;\frac{da_2}{dt}=\bigl(i\Dom-\tfrac{\gamma_2}{2}\bigr)a_2+i\beta\,a_1^2\;}
  \label{eq:a2}\\[4pt]
  \boxed{\;\frac{da_1}{dt}=-\tfrac{\gamma_1}{2}a_1
                        +i\beta^*\,a_1^*\,a_2
                        +i\alpha_3\,|a_1|^2 a_1
                        +\sqrt{\gamma_1^{(\text{rad})}}\,s_+\;}
  \label{eq:a1}
\end{align}
where $\Dom\equiv 2\omega_1-\omega_2$, $s_+$ is the input wave amplitude
normalised so that $|s_+|^2$ is incident power, and the input--output
relation is $s_-=-s_++\sqrt{\gamma_1^{(\text{rad})}}\,a_1$.  The total
FM decay is $\gamma_1=\gamma_1^{(\text{rad})}+\gamma_1^{(\text{nr})}$;
for the SH mode we have, by design,
$\gamma_2\simeq\gamma_2^{(\text{nr})}$.

\textbf{The SHG and DFG couplings carry the same coupling constant
$\beta$.}  This is the Manley--Rowe / energy-conservation requirement and
can be checked directly: in the lossless limit ($\gamma_n,s_+\!\to 0$, no
Kerr), differentiating $|a_n|^{2}$ from the bare TCMT equations
\eqref{eq:a2}--\eqref{eq:a1} gives, per channel,
\begin{equation}
  \frac{d|a_2|^{2}}{dt}=2\,\mathrm{Im}\bigl[\beta\,a_1^{*2}a_2\bigr],
  \qquad
  \frac{d|a_1|^{2}}{dt}=-2\,\mathrm{Im}\bigl[\beta^{(1)}\,a_1^{*2}a_2\bigr],
  \label{eq:MR-perchannel}
\end{equation}
where $\beta^{(1)}$ denotes the (a priori unknown) coupling appearing
in the FM equation \eqref{eq:a1} (we wrote $i\beta^{*}\,a_1^{*}a_2$
there anticipating the result $\beta^{(1)}=\beta$).  The total
mode-energy rate is therefore
\begin{equation}
  \frac{d}{dt}(|a_1|^{2}+|a_2|^{2})
  =2\,\mathrm{Im}\!\bigl[(\beta-\beta^{(1)})\,a_1^{*2}a_2\bigr],
  \label{eq:MR-total}
\end{equation}
which vanishes for arbitrary $a_1,a_2$ if and only if
$\beta^{(1)}=\beta$, identifying the SH-equation coupling and the
FM-equation coupling as complex conjugates of a single $\beta$.  The
same conclusion follows from the Hamiltonian formulation in
\cite[Eqs.~8--9]{Rodriguez2007}, where energy conservation is built
in from the start.  \textbf{The factor $2$ in the DFG polarisation
\eqref{eq:PDFG} is exactly what produces this equality} -- without
it the SHG and DFG couplings would differ by a factor of $2$ and
energy would not be conserved.  The original note's
``$\kappa$ vs.\ $\kappa/2$'' form violates Manley--Rowe.

\paragraph{Derivation of the modal coefficients $\beta$ and $\alpha_3$.}
The boxed equations \eqref{eq:a2}--\eqref{eq:a1} are obtained by
projecting Maxwell's equations onto $\f_n^*$, starting from the wave
equation \eqref{eq:waveeq-SH}.  We sketch the steps for $\beta$ here;
$\alpha_3$ follows identically with the appropriate $\chithree$ source.

Substitute the single-mode ansatz $\mathbf E(\ri,2\omega_1)=a_2(t)\f_2(\ri)$
into \eqref{eq:waveeq-SH}.  In the slow-envelope limit the
$\partial_t^{2}\mathbf E$ piece of the wave operator contributes
$-2i\omega_2\,\dot a_2\,\f_2$; the spatial part
$[\nabla\!\times\!\nabla\!\times\,-\,k_2^{2}\eps_r]\f_2$ vanishes
on-resonance by the modal definition, and produces $2(\omega_2-2\omega_1)
\omega_2 \eps_0 \tilde\eps_2 \f_2/c^{2}=-2\omega_2\Dom \eps_0\tilde\eps_2\f_2/c^{2}$
off-resonance.  Multiplying through by $\f_2^*(\ri)$, integrating
over space, and using the Brillouin energy normalisation
\eqref{eq:Brillouin} (so that $\int \eps_0 \tilde\eps_2 |\f_2|^{2}\,\dthreer
=1$ in our units), the LHS reduces to
\begin{equation*}
  \mathrm{LHS}\;=\;\bigl[\,-2i\omega_2\,\dot a_2 + 2\omega_2\Dom\,a_2\bigr]\cdot
                   (1/\eps_0).
\end{equation*}
The RHS, using $E_j(\omega_1)E_k(\omega_1)=a_1^{2}f_{1,j}f_{1,k}$ in
\eqref{eq:PSHG}, becomes
\begin{equation*}
  \mathrm{RHS}\;=\;k_2^{2}\,a_1^{2}\!\int\!\chitwo_{ijk}\,f_{2,i}^{*}\,f_{1,j}\,f_{1,k}\,\dthreer
            \;=\;k_2^{2}\,a_1^{2}\,M.
\end{equation*}
Equating and dividing by the prefactor of $\dot a_2$ identifies the
free-evolution term $i\Dom\,a_2$ already present in \eqref{eq:a2} and
the SHG drive $i\beta\,a_1^{2}$ with
\begin{equation*}
  \beta = \frac{\eps_0\,k_2^{2}\,c^{2}}{4\,\omega_2}\,M
        = \frac{\eps_0\,\omega_2}{4}\,M
        = \frac{\eps_0\,\omega_1}{2}\,M
\end{equation*}
(using $k_2=\omega_2/c$ and $\omega_2=2\omega_1$).  The conventional
prefactor written in \eqref{eq:beta} below uses $\eps_0\omega_1/4$
rather than $\eps_0\omega_1/2$; the precise numerical coefficient
depends on whether one absorbs $\f_2$-normalisation conventions
$\eps_0\tilde\eps_n |\f_n|^2\to 2$ or $\to 1$ in the Brillouin energy
integral, and on factor-of-2 ambiguities between $|E|^{2}$ and
$|\f_n|^2 |a_n|^2$ when complex envelopes are extracted as in
\eqref{eq:PSHG}--\eqref{eq:PDFG}.  We adopt the standard cavity-QED
choice~\cite{Rodriguez2007,SuhWangFan2004} giving
\begin{equation}
  \boxed{\;\beta=\frac{\eps_0\,\omega_1}{4}\!\int\!\chitwo_{ijk}(\ri)\,
                f_{2,i}^*(\ri)\,f_{1,j}(\ri)\,f_{1,k}(\ri)\,\dthreer\;}
  \qquad[\text{units:}\ \mathrm{J}^{-1/2}\mathrm{s}^{-1}].
\label{eq:beta}
\end{equation}
The same projection applied to the wave equation at $\omega_1$ with
the SPM source $P^{(3)}(\omega_1)=3\eps_0\chithree|E|^{2}E$ produces
the FM equation \eqref{eq:a1} with the bare Kerr coefficient
\begin{equation}
  \boxed{\;\alpha_3=\frac{3\,\eps_0\,\omega_1}{8}\!\int\!\chithree_{ijkl}(\ri)\,
                  f_{1,i}^*\,f_{1,j}\,f_{1,k}\,f_{1,l}^*\,\dthreer\;}
  \qquad[\text{units:}\ \mathrm{J}^{-1}\mathrm{s}^{-1}].
\label{eq:alpha3}
\end{equation}
The factor $3$ is the SPM Boyd degeneracy of
\S\ref{sec:conventions}; the factor $1/8$ tracks the same
$\eps_0\omega_1/4$ envelope/normalisation chain as $\beta$ in
\eqref{eq:beta}, with one additional factor of $1/2$ from extracting
the $|E|^{2}E$ piece of $|\mathcal{E}|^{2}\mathcal{E}$ at frequency
$\omega_1$ in the Boyd field convention (\S\ref{sec:conventions}).
The full derivation parallels the $\beta$ derivation above; see
\cite[Eq.~(8)]{Rodriguez2007} for an explicit treatment.
The DFG source $P^{(2)}(\omega_1)$ of \eqref{eq:PDFG} projects onto
$\f_1^*$ to give the term $i\beta^{*}a_1^{*}a_2$ in \eqref{eq:a1};
this is most cleanly seen via the conservative Hamiltonian
$H=\beta^{*}a_1^{*2}a_2+\beta\,a_1^{2}a_2^{*}$, whose equations of
motion $i\dot a_n=\partial H/\partial a_n^{*}$ produce \eqref{eq:a2}
and the DFG term in \eqref{eq:a1} simultaneously with the
\emph{same} $\beta$, automatically satisfying \eqref{eq:MR-total}.

The numerical prefactors $\eps_0\omega_1/4$ and $3\eps_0\omega_1/8$
depend on the chosen normalisation and field convention; the
structure of the integrals -- full contraction of the tensor with the
vectorial mode profiles, weighted by $\eps_0$ from the
polarisation--field relation -- is convention-independent and is the
practical input for full-wave numerical evaluation.  The factor $3$
in $\alpha_3$ is the SPM degeneracy from
Section~\ref{sec:conventions} and was absent in the original note.

\paragraph{Equivalence with the polarisation-tensor derivation.}
The TCMT route here and the Green's-function route of
Section~\ref{sec:tensorder} are the same calculation in different
representations.  Substituting the QNM Green's function
\eqref{eq:G-QNM} into \eqref{eq:Esh-Green} gives
$\mathbf E(\ri,2\omega_1)=a_2\,\f_2(\ri)$ with the steady-state
$a_2\propto a_1^{2}/(\Dom+i\gamma_2/2)$ of \eqref{eq:a2}, with
overall prefactor fixed by \eqref{eq:beta}.  This consistency
determines the QNM Green's-function normalisation in \eqref{eq:G-QNM}
and confirms that the cascaded $\chithree$ identified there and
in the adiabatic-elimination route of Section~\ref{sec:adiab}
(eq.~\eqref{eq:chi3casc}) agree.

% =======================================================================
\section{Adiabatic elimination and the effective $\chithree$}\label{sec:adiab}

Assume $|\Dom|\gg\gamma_2/2$ and $|\dot a_2|\ll|\Dom\,a_2|$ (slow
modulation of the FM compared with the SH response time).  Setting
$da_2/dt\to 0$ in \eqref{eq:a2},
\begin{equation}
  a_2 = \frac{i\beta\,a_1^2}{-i\Dom+\gamma_2/2}
      = \frac{-\beta\,a_1^2}{\Dom+i\gamma_2/2}
      \simeq \frac{-\beta\,a_1^2}{\Dom}\!\left(1-\frac{i\gamma_2}{2\Dom}\right).
\end{equation}
Substituting into \eqref{eq:a1},
\begin{equation}
  \frac{da_1}{dt}=
   \Bigl[-\tfrac{\gamma_1}{2}-\underbrace{\tfrac{|\beta|^2\gamma_2}{2\Dom^{2}}}_{\text{cascaded loss}}|a_1|^2\Bigr]a_1
   +i\Bigl[\underbrace{\alpha_3-\tfrac{|\beta|^2}{\Dom}}_{\text{effective Kerr}}\Bigr]|a_1|^2 a_1
   +\sqrt{\gamma_1^{(\text{rad})}}\,s_+ .
\end{equation}

\paragraph{(i) Effective Kerr coefficient.}
\begin{equation}
  \boxed{\;\alpha_3^{\text{eff}}=\alpha_3-\frac{|\beta|^2}{\Dom}\;}
\end{equation}
The cascaded contribution can have \emph{either} sign relative to the
bare Kerr response, depending on whether $\Dom\equiv 2\omega_1-\omega_2$
is positive or negative.  This sign-tunability is one of the hallmarks
of cascaded nonlinearity: the same metasurface can be tuned through
zero, used for self-defocusing, or used for self-focusing simply by
walking the SH resonance through twice the FM frequency.  This was used
dramatically in soliton compression experiments \cite{DeSalvo1992,Liu1999}.

\paragraph{(ii) Two-photon-like loss.}
The cascaded loss coefficient -- defined as the coefficient $X$ in the
bracket form $da_1/dt = [-\gamma_1/2 - X|a_1|^{2}]a_1 + \ldots$ on
\eqref{eq:Pcasc-nonlocal} above -- is, in the adiabatic limit
$|\Dom|\gg\gamma_2/2$,
\begin{equation}
  \gamma_3^\text{eff}\;=\;\frac{|\beta|^{2}\,\gamma_2}{2\,\Dom^{2}},
  \label{eq:gamma3-adiab}
\end{equation}
with full Lorentzian form
$\gamma_3^\text{eff}(\Dom)=|\beta|^{2}(\gamma_2/2)/[\Dom^{2}+(\gamma_2/2)^{2}]$
peaking at $\Dom=0$ with value $2|\beta|^{2}/\gamma_2$.  It is the
analogue of two-photon absorption (here governed by $\gamma_2$ rather
than by $\mathrm{Im}\,\chithree$) and is unavoidable as the price of
the resonant enhancement.  In the same convention, the cascaded
loss-to-Kerr ratio at the design point is
$\gamma_3^\text{eff}/|\alpha_3^\text{eff}|=\gamma_2/(2|\Dom|)$.

\bigskip

To compare with bulk material parameters, invert the definition of
$\alpha_3$ in \eqref{eq:alpha3}: a uniform bulk $\chithree_{\text{eff}}$
filling the mode produces
$\alpha_3=(3\eps_0\omega_1/8)\,\chithree_{\text{eff}}\!\int\!|\f_1|^4\dthreer$,
so the cascaded contribution corresponds to
\begin{equation}
  \chithree_{\text{casc}}=-\frac{|\beta|^2/\Dom}{(3\eps_0\omega_1/8)\!\int\!|\f_1|^4\dthreer}
                       =-\frac{8}{3\eps_0\omega_1\Dom}\,
                            \frac{|\beta|^2}{\int|\f_1|^4\dthreer}.
\end{equation}
Inserting the explicit $\beta=\eps_0\omega_1\,I/4$ from \eqref{eq:beta} (with
$I\equiv\int\chitwo_{ijk}\,f_{2,i}^*f_{1,j}f_{1,k}\dthreer$) and using
$|\beta|^2=\eps_0^2\omega_1^2|I|^2/16$,
\begin{equation}
  \boxed{\;\chithree_{\text{casc}}=-\frac{\eps_0\,\omega_1}{6\,\Dom}\,
       \frac{\bigl|\!\int\!\chitwo_{ijk}(\ri)\,f_{2,i}^*\,f_{1,j}\,f_{1,k}\dthreer\bigr|^{2}}
            {\int|\f_1|^4\dthreer}\;.}
\label{eq:chi3casc}
\end{equation}
\paragraph{Explicit $n$-dependence.}
The $n$ that one might have expected in a ``naïve'' (plane-wave-style)
derivation is hidden in the energy normalisation
\eqref{eq:Brillouin}: at unit stored energy $|a_n|^2=1$~J, the typical
field magnitude in the dielectric is
\begin{equation}
  |\f_n(\ri)|\;\sim\;\frac{1}{\sqrt{\eps_0\,n_n^{\,2}\,V_\text{eff}}},
\end{equation}
where $n_n$ is the effective refractive index seen by mode $n$ (close
to the bulk index at $\omega_n$ for a tightly confined dielectric mode)
and $V_\text{eff}$ is the mode volume.  Substituting these scalings,
\begin{align}
  \int|\f_1|^4\dthreer&\sim\frac{\zeta_1}{\eps_0^{\,2}\,n_1^{\,4}\,V_\text{eff}},
  \\[4pt]
  \Bigl|\!\int\!\chitwo_{ijk}f_{2,i}^*f_{1,j}f_{1,k}\dthreer\Bigr|^{2}
        &\sim\frac{|\chitwo|^{2}\,\widetilde\eta^{2}}{\eps_0^{\,3}\,n_2^{\,2}\,n_1^{\,4}\,V_\text{eff}},
\end{align}
where $\zeta_1,\widetilde\eta$ are purely geometric, dimensionless numbers of
order unity (perfect overlap $\widetilde\eta\to 1$).  All four powers of
$n_1$ cancel between numerator and denominator, and one $\eps_0$
cancels against the $\eps_0$ prefactor in \eqref{eq:chi3casc}, giving
the final explicit form
\begin{equation}
  \boxed{\;\chithree_{\text{casc}}\;\sim\;-\frac{\omega_1}{6\,n_2^{\,2}\,\Dom}\;|\chitwo|^{2}\;\eta_\text{geom}\;,
        \qquad \eta_\text{geom}\equiv\frac{\widetilde\eta^{\,2}}{\zeta_1}\sim 0.01\!-\!0.5\;}
\label{eq:chi3casc_explicit}
\end{equation}
with $\eta_\text{geom}$ a purely geometric overlap.  Three things are
worth emphasising about this form:
\begin{itemize}
  \item \textbf{Only} $n_2$ \textbf{(the SH-mode effective index) survives.}
        The $n_1$ powers cancel between numerator and denominator
        because $|\f_1|^4$ appears symmetrically in both, leaving the
        SH-mode normalisation $|\f_2|^2\propto 1/(\eps_0 n_2^{\,2}V)$
        as the only residue.  Physically, modes in high-index material
        concentrate \emph{less} field per unit stored energy, so the
        $\f_2{:}\f_1\f_1$ coupling is weaker by $1/n_2$.
  \item \textbf{The dielectric constant $\eps_0$ does not appear} in the
        final result.  This is a non-trivial consistency check: $\chi^{(3)}_\text{eff}$
        is a material parameter (m$^2$/V$^2$ in SI), and dimensional
        analysis of $|\chitwo|^{2}$ alone gives the right units.
  \item The overall coefficient (here $1/6$) is convention-dependent;
        the structure $\omega_1/(n_2^{\,2}\Dom)\cdot|\chitwo|^{2}\cdot\eta_\text{geom}$
        is universal.
\end{itemize}

% =======================================================================
\section{The bandwidth--enhancement sum rule}\label{sec:sumrule}

The peak enhancement is bounded above by the SH linewidth.  The
cleanest derivation is from the AC frequency response of the cascade.
If $|a_1|^{2}$ is modulated at frequency $\Omega$ (so the FM intensity
has a Fourier component $\delta|a_1|^{2}\,e^{-i\Omega t}$ on top of a
DC value), the SH equation \eqref{eq:a2} generates a corresponding
SH-amplitude sideband
$\delta a_2 = i\beta\,\delta(a_1^{2})/(-\!i(\Dom-\Omega)+\gamma_2/2)$,
and the DFG back-action on the FM acts as an effective \emph{frequency-dependent}
Kerr coefficient
\begin{equation}
  \alpha_3^{\text{eff}}(\Omega)
   \;=\;\frac{i\,|\beta|^{2}}{\gamma_2/2+i(\Omega-\Dom)},
  \qquad
  |\alpha_3^{\text{eff}}(\Omega)|
   \;=\;\frac{|\beta|^{2}}{\sqrt{(\gamma_2/2)^{2}+(\Omega-\Dom)^{2}}}.
  \label{eq:alpha3-AC}
\end{equation}
The \emph{Lorentzian} in $\Omega$ here is $|\alpha_3^{\text{eff}}(\Omega)|^{2}$,
centred at $\Omega=\Dom$ with FWHM $=\gamma_2$, peak
$(2|\beta|^{2}/\gamma_2)^{2}=4|\beta|^{4}/\gamma_2^{2}$, and area
\begin{equation}
  \int\!d\Omega\,|\alpha_3^{\text{eff}}(\Omega)|^{2}
   \;=\;\frac{2\pi\,|\beta|^{4}}{\gamma_2}.
  \label{eq:area-rule}
\end{equation}
$|\alpha_3^{\text{eff}}(\Omega)|$ itself is the square root of this
Lorentzian, so it is \emph{broader}: FWHM$(|\alpha|)=\sqrt{3}\,\gamma_2$
(taking the square root widens by $\sqrt{3}$).  Its peak is
$|\alpha_3^{\text{eff}}|_\text{peak}=|\beta|^{2}/(\gamma_2/2)=2|\beta|^{2}/\gamma_2$
at $\Omega=\Dom$, and the peak-times-FWHM product
\begin{equation}
  |\alpha_3^{\text{eff}}|_\text{peak}\times \text{FWHM}(|\alpha_3^{\text{eff}}|)
  \;=\;2\sqrt{3}\,|\beta|^{2},
\end{equation}
is conserved as $Q_2$ varies (since $|\beta|^{2}\propto|\chitwo|^{2}\eta_\text{geom}/n_2^{\,2}$
of Section~\ref{sec:adiab} is set by material constants and overlap
geometry alone, independent of $Q_2$).  The cascaded coefficient has
resonant gain when the FM modulation frequency matches the SH detuning.

In design terms, at the conventional operating point $|\Dom|\!\sim\!\gamma_2$:
$|\alpha_3^{\text{eff}}|_\text{max}\!\sim\!|\beta|^{2}/\gamma_2$ and
the usable bandwidth is $\Delta\Omega\!\sim\!\gamma_2$, so combining
with \eqref{eq:chi3casc_explicit} (using $\omega_1/\gamma_2 = Q_2/2$
since $\gamma_2 = \omega_2/Q_2 = 2\omega_1/Q_2$):
\begin{equation}
  \bigl|\chithree_\text{casc}\bigr|_\text{max}\;\sim\;
       \frac{\omega_1}{6\,n_2^{\,2}\,\gamma_2}\,|\chitwo|^2\,\eta_\text{geom}
       \;=\;\frac{Q_2}{12\,n_2^{\,2}}\,|\chitwo|^2\,\eta_\text{geom}\,,
  \label{eq:chi3-max}
\end{equation}
and (multiplying by $\Delta f_\text{usable}\sim\gamma_2/(2\pi)$, so that
$\gamma_2$ cancels)
\begin{equation}
  \boxed{\;\bigl|\chithree_\text{casc}\bigr|_\text{max}\,\times\,\Delta f_\text{usable}
          \;\lesssim\;\frac{\omega_1\,|\chitwo|^2\,\eta_\text{geom}}{12\pi\,n_2^{\,2}}
          \quad\text{(independent of }Q_2\text{).}\;}
  \label{eq:sum-rule}
\end{equation}
$Q_2$ does not enter the product: \textbf{resonant enhancement
redistributes nonlinear action into a tall, narrow peak; it does not
increase the area under the curve.}  This is the same rule that
limits resonant linear absorption (oscillator-strength sum rule),
purely classical, and it is what disqualifies high-$Q$
double-resonance schemes from delivering ``fast'' $\chithree$ in the
ultrafast sense.  For the nominal numbers below ($Q_2\sim 4{\times}10^3$
at $\lambda_1=1.55\,\mu$m, so $\gamma_2=\omega_2/Q_2=2\omega_1/Q_2$),
$\Delta f_\text{usable}\sim\gamma_2/(2\pi)\simeq 97$~GHz: narrowband
by Kerr-electronics standards (THz), broadband by photonic-cavity
standards.

For non-resonant (phase-mismatched bulk) cascading the same sum rule
applies with $\gamma_2$ replaced by the inverse coherence time
$|\Delta k|\,v_g$, giving modest enhancement over a much larger
bandwidth -- the regime exploited in soliton compression and bulk
all-optical switching.

Equation \eqref{eq:alpha3-AC} describes the cascade response to a
modulation of $|a_1|^{2}$ at frequency $\Omega$ (relevant for slow
envelope modulation of the FM amplitude); the input is the pump
intensity and the AC-peak/DC enhancement of $|\alpha_3^\text{eff}|$ is
$|\Dom|/(\gamma_2/2)$ asymptotically.  A \emph{distinct} transfer
function (different input, different output) is treated in
Section~\ref{sec:applications}, item~6: there one modulates the
SH-cavity detuning $\Dom(t)$ at frequency $\Omega_\text{mod}$ -- this
is the cascade-based amplitude modulator -- and its AC-peak/DC
enhancement is $\sqrt{\Dom_0^{2}+(\gamma_2/2)^{2}}/\gamma_2\approx
|\Dom_0|/\gamma_2$ (eq.~\eqref{eq:AC-peak-enhancement}), smaller by a
factor of $2$ than the FM-modulation result because the SH-detuning
modulation only drives one sideband per Fourier component of $\Dom(t)$.

% =======================================================================
\section{Tensor structure and phase matching}\label{sec:tensor}

The overlap that defines $\beta$,
\begin{equation*}
  \beta\propto\!\int\!\chitwo_{ijk}(\ri)\,f_{2,i}^*\,f_{1,j}\,f_{1,k}\dthreer,
\end{equation*}
vanishes identically for several reasons, each of which must be
addressed in the metasurface design:
\begin{enumerate}
  \item \textbf{Bulk inversion symmetry.} If a meta-atom has inversion
        symmetry, $\chitwo_{ijk}\equiv 0$.  Hence one needs a
        non-centrosymmetric crystal substrate (LiNbO$_3$, point group
        $3m$; GaAs/GaP, point group $\bar{4}3m$ ($T_d$); AlGaAs;\,\dots)
        \emph{and} a meta-atom geometry that does not impose inversion
        symmetry on the mode profile.

  \item \textbf{Mode-symmetry mismatch.} Even with $\chitwo_{ijk}\neq 0$
        in the material, the integrand has a definite parity under any
        spatial symmetry shared by both modes.  If the FM mode is, say,
        magnetic-dipole-like (odd under a mirror) and the SH is
        electric-dipole-like (even), the integrand may be odd under
        that mirror and the integral vanishes.  Quasi-BIC engineering
        is precisely the controlled breaking of one such symmetry to
        lift this constraint while keeping $Q_2$ large.

  \item \textbf{Quasi-phase matching across the metasurface period.}
        For an extended metasurface, the integrand carries Bloch phase
        factors $e^{i(2\bvec{k}_1-\bvec{k}_2)\cdot\bvec{R}}$.  Periodic
        poling or spatial modulation of $\chitwo_{ijk}(\ri)$ across
        unit cells provides a reciprocal-lattice vector $\bvec{G}$ that
        compensates the residual momentum
        $2\bvec{k}_1-\bvec{k}_2-\bvec{G}=0$.  In a single-unit-cell
        metasurface with both modes at $\bvec{k}_\parallel=0$
        ($\Gamma$ point), this momentum-matching is automatic and the
        only requirement is on the \emph{internal} symmetry of the mode
        profiles.
\end{enumerate}

For LiNbO$_3$ (point group $3m$), the non-zero independent components of
$\chitwo_{ijk}$ in contracted notation are
$d_{15}=d_{31}$, $d_{22}=-d_{21}=-d_{16}$, and $d_{33}$, with
$|d_{33}|\simeq 25$~pm/V at 1064~nm being the largest.  With the
$\hat{z}$ axis parallel to the optic axis, the largest overlap is
achieved by mode profiles with strong $|f_{1,z}|^{2}f_{2,z}^{*}$ content
over the meta-atom volume.  Z-cut LiNbO$_3$ with field components
engineered along the optic axis is therefore preferred.  The same
considerations apply, with the appropriate tensor structure, for
$\bar{4}3m$ semiconductors (largest component $d_{14}=d_{25}=d_{36}\simeq 100$~pm/V
for GaAs).

For LiNbO$_3$, $\chitwo\equiv 2d_{33}\simeq 50$~pm/V $=5\times10^{-11}$~m/V;
for GaAs, $\chitwo\simeq 2\times10^{-10}$~m/V.

% =======================================================================
\section{Metasurface implementation: BIC and quasi-BIC platforms}

The high-$Q_2$ requirement is the key technical hurdle, and bound states
in the continuum (BIC) provide a particularly elegant route
\cite{Hsu2016}.  A symmetry-protected BIC is a guided-wave-like solution
that lies inside the radiation continuum but is \emph{symmetry-forbidden}
from coupling to plane waves at $\bvec{k}_\parallel=0$.  Realistic
structures are quasi-BICs (qBICs): finite arrays, slight disorder, or
deliberate symmetry breaking convert the ideal BIC into a high-$Q$ leaky
resonance, whose $Q$ scales as $1/\alpha^{2}$ with the asymmetry
parameter $\alpha$ \cite{Koshelev2018,Bulgakov2017}.  Demonstrated $Q$ factors in
dielectric metasurfaces are routinely $10^{3}$--$10^{4}$ and have
reached $\sim 10^{5}$ in carefully optimised devices.

Two practical architectures:
\begin{enumerate}
  \item[\textbf{(a)}] \textbf{Two-mode qBIC pair.}  Both the FM and SH
        are qBICs of the same array, with the SH qBIC very close to the
        symmetry-protected limit (large $Q_2$) and the FM qBIC
        moderately broadened ($Q_1\sim 10^{2}$--$10^{3}$) to allow
        efficient in-coupling at $\omega_1$.  The SH leakage to the
        radiation continuum is irrelevant for the cascaded process --
        we \emph{don't want} SH out-coupling here.

  \item[\textbf{(b)}] \textbf{FM Mie--qBIC + SH symmetry-protected BIC.}
        The FM is a low-$Q$ Mie resonance (dipolar, easy to pump); the
        SH lives in a true BIC mode of the array, with $\gamma_2$ set
        by absorption only.  Best for maximum enhancement, hardest to
        fabricate.
\end{enumerate}

In both cases the requirement on the SH mode is purely \emph{internal}:
high $Q_2$, large field overlap with the FM-squared, no need for any
external connection.

A subtle point: the cascaded $\chithree$ effect manifests at the FM, so
the experimental signature is \textbf{intensity-dependent
reflection/transmission at $\omega_1$} (self-phase modulation,
intensity-dependent line shape, optical bistability for sufficient pump),
not SH light leaving the sample.  In fact, a \emph{good} implementation
produces \emph{no measurable} SH radiation while the back-action on the
fundamental is maximum.  This is also the experimental fingerprint that
distinguishes the cascaded process from the bare Kerr response.

% =======================================================================
\section{A worked numerical estimate (LiNbO$_3$ metasurface)}\label{sec:numerical}

Take

\begin{center}
\begin{tabular}{ll}
\toprule
quantity & value \\
\midrule
$\lambda_1$               & $1.55\,\mu$m \\
$\omega_1$                & $1.22\times10^{15}$ rad/s \\
$n_1,\,n_2$               & $2.21,\,2.26$ \\
$n_2^{\,2}$               & $\simeq 5.1$ \\
$|\chitwo|$ ($=2d_{33}$)  & $5\times10^{-11}$ m/V \\
$|\chithree|_\text{native}$ & $\sim 2\times10^{-21}$ m$^2$/V$^2$ \\
$Q_2$                     & $4\times10^{3}$ \\
$\gamma_2=\omega_2/Q_2=2\omega_1/Q_2$ & $6.1\times10^{11}$ rad/s
                                       ($\gamma_2/2\pi\simeq 97$~GHz) \\
$\Dom$ (safely adiabatic, $|\Dom|\gtrsim 2\gamma_2$) &
   $2\gamma_2 = 4\omega_1/Q_2 = 1.22\times10^{12}$ rad/s \\
$\eta_\text{geom}$ (purely geometric overlap, eq.~\eqref{eq:chi3casc_explicit})
                          & $\sim 0.3$ \\
\bottomrule
\end{tabular}
\end{center}

Using the explicit form \eqref{eq:chi3casc_explicit} with $\Dom=2\gamma_2$,
so $\omega_1/\Dom = Q_2/4$:
\begin{align*}
  \bigl|\chithree_\text{casc}\bigr|\;&\sim\;
        \frac{\omega_1}{6\,n_2^{\,2}\,\Dom}\,|\chitwo|^{2}\,\eta_\text{geom}
   \;=\;\frac{Q_2}{24\,n_2^{\,2}}\;|\chitwo|^{2}\,\eta_\text{geom}\\[2pt]
   &\sim\;\frac{4000}{24\times 5.1}\times 0.3\times|\chitwo|^{2}
   \;\sim\;10\times|\chitwo|^{2},
\end{align*}
that is
\begin{equation*}
  \bigl|\chithree_\text{casc}\bigr|\sim
       10\times(5\times10^{-11})^{2}\ \mathrm{m^{2}/V^{2}}
       \;\simeq\;2.5\times10^{-20}\ \mathrm{m^{2}/V^{2}},
\end{equation*}
i.e.\ roughly \textbf{12$\boldsymbol{\times}$ the native bulk}
$\chithree$ of LiNbO$_3$, available over a usable bandwidth
$\sim\gamma_2/(2\pi)\simeq 97$~GHz centred at the FM resonance
(Section~\ref{sec:sumrule}).  The $1/n_2^{\,2}$ factor costs a factor
of $\sim 5$ relative to a ``naïve'' estimate that drops it -- it is
not negligible, and an honest design has to compensate it with either
larger $\eta_\text{geom}$ or larger $Q_2/|\Dom|$.

\paragraph{Operating-point sensitivity.}
Pushing $\Dom$ from the safely adiabatic $2\gamma_2$ down toward the
edge-of-validity $\Dom\!\sim\!\gamma_2$ (Section~\ref{sec:validity},
item~1) roughly doubles the enhancement to $R\!\sim\!25$, at the
cost of working at the edge of the adiabatic approximation where the
cascaded TPA (Section~\ref{sec:adiab}(ii)) becomes comparable to the
Kerr response: indeed at $|\Dom|=\gamma_2$, the loss-to-Kerr ratio
$\gamma_\text{eff}^{(2)}/|\alpha_3^\text{eff,casc}|=\gamma_2/(2|\Dom|)=1/2$.
Several earlier estimates in the cascade literature implicitly sit at
this operating point and quote $R\!\sim\!25$ for LiNbO$_3$.  We adopt
the conservative $\Dom=2\gamma_2$, $R\!\sim\!12$ throughout the rest of
this note, and remind the reader of the factor-of-$\sim\!2$ headroom
when interpreting numerical claims downstream.

\paragraph{Prefactor uncertainty.}  The numerical estimate
$R\!\sim\!12$ carries an additional $O(1)$ multiplicative uncertainty
inherited from the convention dependence of the modal-coupling prefactor
in $\beta=(\eps_0\omega_1/4)\,M$ (see the discussion before
\eqref{eq:beta}: the alternative $\eps_0\omega_1/2$ convention would
shift $|\beta|^2$, and hence $R$, by a factor of $4$).  Together with
the $\eta_\text{geom}\!\sim\!0.3$ assumption -- which is itself
geometry-specific -- the $R$ values quoted here should be read as
\emph{order-of-magnitude} predictions; settling the prefactor requires
a full-wave evaluation of the modes and overlap integrals in a
specific meta-atom design.

With $Q_2=10^{4}$ (achievable in optimised qBIC LiNbO$_3$
metasurfaces) and the same $\eta_\text{geom}=0.3$, the estimate gives
$\sim 30\times$ enhancement at the cost of a $\sim 40$~GHz bandwidth.

For comparison, a phase-mismatched bulk PPLN slab at the same
$|\chitwo|$ gives non-resonant cascaded $\chithree$ via the
SVEA-derived formula \eqref{eq:chi3casc-bulk}.  Inserting
$|\Delta k|\sim 3\!\times\!10^{5}$~m$^{-1}$ (corresponding to a
LiNbO$_3$ coherence length $\pi/|\Delta k|\sim 10\,\mu$m at telecom;
the bare $\Delta n\!\simeq\!0.05$ between $\omega_1$ and $2\omega_1$
gives $|\Delta k|\!\simeq\!4\!\times\!10^{5}$~m$^{-1}$, so this is
a representative order-of-magnitude figure), $n_2\simeq 2.26$, and
$\omega_1/c\simeq 4\!\times\!10^{6}$~m$^{-1}$,
\begin{equation*}
  \bigl|\chithree_\text{casc,\,bulk}\bigr|\;\sim\;
    \frac{2\cdot 4\!\times\!10^{6}}{3\cdot 2.26\cdot 3\!\times\!10^{5}}
    \cdot (5\!\times\!10^{-11})^{2}
  \;\sim\;10^{-20}\ \mathrm{m^{2}/V^{2}},
\end{equation*}
roughly $5\times$ the native $\chithree_\mathrm{LiNbO_3}$, available
over $\sim\!10$~THz bandwidth -- the regime exploited in soliton
compression and bulk all-optical switching.

Comparing $|\chithree|\!\times\!\Delta f$ products honestly: bulk
PPLN gives $|\chithree_\text{casc,\,bulk}|\!\times\!\Delta f\sim
10^{-20}\cdot 10^{13}\!=\!10^{-7}$~m$^{2}$V$^{-2}$Hz;
the metasurface at $Q_2=4{\times}10^{3}$ gives
$|\chithree_\text{casc,\,meta}|_\text{peak}\!\times\!\Delta f\sim
2.5\!\times\!10^{-20}\cdot 10^{11}\!=\!2.5\!\times\!10^{-9}$~m$^{2}$V$^{-2}$Hz,
about a factor $\sim\!40$ \emph{less} than bulk (the $\eta_\text{geom}$
and $1/n_2^{2}$ penalties are not free).  The metasurface's appeal
is therefore not in the integrated nonlinear strength but in
\textbf{the chip-scale form factor and CW-low-threshold operation}:
narrow-band high-peak operation in a $\sim 100$-nm-thick film, with
$\mu$W-scale CW pumping rather than the $>\!100$~W pulsed pumping
of bulk PPLN cascading.  Both regimes have their applications; the
high-$Q$ regime is suited to CW-pumped bistability, parametric
oscillation thresholds and frequency-comb stabilisation, rather
than to ultrafast switching.

% =======================================================================
\section{Comparison with a native-$\chithree$ resonance: when does cascading help?}\label{sec:comparison}

The numerical estimate of the previous section -- $|\chithree_\text{casc}|
\simeq 12\,|\chithree_\text{native}|$ at $Q_2=4{\times}10^{3}$ and
$\Dom=2\gamma_2$ (Section~\ref{sec:numerical}) -- begs
the question: is the cascaded route actually preferable to engineering
a single high-$Q$ resonance that enhances the \emph{native} bulk
$\chithree$?  Both architectures use the same trick (a high-$Q$
dielectric resonance at the fundamental) to build up intracavity
intensity; they differ only in what mechanism inside the cavity then
turns that intensity into a phase shift.  The comparison is delicate
because the two architectures have different bandwidth structures and
different saturation channels, and the naive ratio
$|\chithree_\text{casc}|/|\chithree_\text{native}|$ does not by itself
settle the matter.  We answer it here within the same TCMT framework.

% --------------------------------------------------------------
\subsection*{Unified TCMT for both architectures}

Consider a single FM cavity of frequency $\omega_1$ and total decay
$\gamma_1$, identical in both architectures.  The TCMT equations
\eqref{eq:a2}--\eqref{eq:a1} already include both candidate sources of
$\chithree_\text{eff}$ on the FM: the bulk-Kerr coefficient $\alpha_3$
and the cascaded back-action via $\beta$ and the SH mode $a_2$.  The
two architectures are limits of the same model:

\begin{itemize}\setlength{\itemsep}{3pt}
  \item \textbf{Architecture A (native-$\chithree$ resonance):}
        $\beta\to 0$ (no SH mode engineered, or $|\Dom|\!\to\!\infty$).
        Only the FM resonance is used; the cubic on the FM is the bulk
        $\alpha_3$.
  \item \textbf{Architecture B (cascaded $\chitwo{:}\chitwo$):}
        $\alpha_3\to 0$ (work in a material with negligible native
        Kerr, or simply isolate the cascaded contribution from a
        full-Kerr model).  All the action comes from $\beta$ and the
        SH mode at $\omega_2\approx 2\omega_1$.
  \item \textbf{Architecture A+B (combined):} both terms present;
        relevant for materials with appreciable $\chitwo$ \emph{and}
        $\chithree$.
\end{itemize}

\noindent
Because the FM cavity is the same, the (linear-cavity, on-resonance)
build-up factor $|a_1|^{2}/|s_+|^{2}=4\gamma_1^\text{rad}/\gamma_1^{2}$
(which equals $4/\gamma_1$ in the lossless / fully-radiative limit
$\gamma_1^\text{rad}=\gamma_1$, and $2/\gamma_1$ at the conventional
critical-coupling point $\gamma_1^\text{rad}=\gamma_1/2$ -- the
formulae below assume the former), and the
interaction time $2/\gamma_1$ are identical.  Any architectural
advantage must therefore come from what multiplies these factors.

% --------------------------------------------------------------
\subsection*{Per-photon nonlinear phase: the $Q$-scaling}

At pump frequency $\omega_p=\omega_1$ and small signal (no saturation,
no bistability) the steady-state intracavity intensity is $|a_1|^{2}=
(4/\gamma_1)\,|s_+|^{2}$, so the accumulated nonlinear phase
$\phi_\text{NL}=\alpha_3^\text{eff}\,|a_1|^{2}\cdot(2/\gamma_1)$ per
input photon takes the simple form
\begin{equation}
  \boxed{\;\frac{\phi_\text{NL}}{|s_+|^{2}}\;=\;\frac{8\,\alpha_3^\text{eff}}{\gamma_1^{2}}\;}
  \label{eq:phi-per-photon}
\end{equation}
where $\alpha_3^\text{eff}$ is whichever Kerr coefficient drives the
FM cavity.  For the two architectures,
\begin{align}
  \text{Architecture A:}\qquad &
   \alpha_3^\text{eff,A}=\alpha_3\;\propto\;\chithree_\text{native},
   \label{eq:alpha-A}\\[3pt]
  \text{Architecture B:}\qquad &
   \alpha_3^\text{eff,B}=-\frac{|\beta|^{2}}{\Dom+i\gamma_2/2}
            \;\propto\;\frac{|\chitwo|^{2}\,\eta_\text{geom}}{n_2^{\,2}\,\Dom}
            \;\propto\;Q_2.
   \label{eq:alpha-B}
\end{align}
The structural difference is condensed in one parameter: the
\emph{cascaded enhancement ratio}
\begin{equation}
  \boxed{\;R\;\equiv\;\frac{|\alpha_3^\text{eff,B}|}{|\alpha_3^\text{eff,A}|}
        \;=\;\frac{|\beta|^{2}}{|\alpha_3|\,|\Dom+i\gamma_2/2|}
        \;\stackrel{\Dom\gg\gamma_2}{\simeq}\;
        \frac{|\chithree_\text{casc}|}{|\chithree_\text{native}|}\;.}
  \label{eq:R-def}
\end{equation}
$R$ is the proper figure of merit because it depends only on the
material constants and the SH-resonance design ($Q_2$, $\Dom$, the
overlap $\eta_\text{geom}$), \emph{not} on the FM cavity properties
$Q_1,\,\gamma_1^\text{rad}$ which both architectures share.  Note that
$R$ as defined uses $|\alpha_3^\text{eff,B}|$ (the modulus, including
$i\gamma_2/2$), while the asymptotic identification with
$|\chithree_\text{casc}|/|\chithree_\text{native}|$ uses the real part
of $\alpha_3^\text{eff,B}$ only; the two agree to within $\sim 3\%$ at
$\Dom=2\gamma_2$ and to $\sim 11\%$ at $\Dom=\gamma_2$.

\paragraph{Q-scaling.}  Combining \eqref{eq:phi-per-photon} with
\eqref{eq:alpha-A}--\eqref{eq:alpha-B}:
\begin{align}
  \phi_\text{NL}^\text{A} &\;\propto\;Q_1^{2}\,\chithree_\text{native},
  \label{eq:phi-A-scaling}\\[2pt]
  \phi_\text{NL}^\text{B} &\;\propto\;Q_1^{2}\,Q_2\,\frac{|\chitwo|^{2}\,\eta_\text{geom}}{n_2^{\,2}}.
  \label{eq:phi-B-scaling}
\end{align}
\emph{Both architectures benefit equally from $Q_1$} (the FM-cavity
build-up scales as $Q_1$ and the interaction time scales as $Q_1$,
giving the $Q_1^{2}$ factor common to both).  The cascade adds an
\emph{extra} $Q_2$ factor: increasing the SH quality factor at fixed
detuning ratio $\Dom/\gamma_2$ pulls more action out of the same
intracavity FM intensity.

% --------------------------------------------------------------
\subsection*{The phase--bandwidth locus}

Pushing either $Q_1$ or $Q_2$ to gain phase costs bandwidth.  The
useful operational bandwidth of each architecture is
\begin{align*}
  \text{Architecture A:}\qquad &\Delta f_\text{A}=\gamma_1/(2\pi)
        =\omega_1/(2\pi Q_1),\\
  \text{Architecture B:}\qquad &\Delta f_\text{B}=\min(\gamma_1,\gamma_2)/(2\pi).
\end{align*}
For architecture B, $\gamma_2$ enters because the cascaded coefficient
itself has finite bandwidth (Section~\ref{sec:sumrule} sum rule); whichever
of $\gamma_1$ or $\gamma_2$ is smaller bottlenecks the device.

Plotting $\phi_\text{NL}/|s_+|^{2}$ versus $\Delta f$ in log--log
reveals three distinct families with three distinct slopes:

\begin{center}
\begin{tabular}{lll}
\toprule
Family & Operating mode & Slope on log--log \\
\midrule
A: native, sweep $Q_1$            &
   $\phi\!\propto\! Q_1^{2}$, $\Delta f\!\propto\!1/Q_1$ &
   $-2$ \\
B$_\text{II}$: cascade, fix $Q_1$, sweep $Q_2$ &
   $\phi\!\propto\!Q_2$, $\Delta f\!\propto\!1/Q_2$ &
   $-1$ (sum-rule) \\
B$_\text{I}$: cascade, fix $Q_2$, sweep $Q_1$ &
   $\phi\!\propto\!Q_1^{2}$, $\Delta f\!\propto\!1/Q_1$ &
   $-2$ \\
\bottomrule
\end{tabular}
\end{center}

The native locus A is a single line; the cascade is a two-parameter
\emph{region} in $(\phi,\Delta f)$ bounded by B$_\text{I}$ and
B$_\text{II}$.  Three qualitative features follow:

\begin{enumerate}
  \item \textbf{Vertical offset.} At a given bandwidth, the cascade
    locus is shifted above the native locus by the structural ratio
    $R$ of \eqref{eq:R-def}.  $R\gg 1$ is the regime in which
    cascading wins outright; $R\lesssim 1$ is the regime in which
    native wins.

  \item \textbf{Slopes differ.} A slope-$-1$ line and a slope-$-2$
    line cross. For $R\gg 1$, the crossover happens at very narrow
    bandwidth -- well past anything experimentally interesting,
    so cascade wins everywhere.  For $R\sim 1$, the crossover sits in
    the experimentally relevant range and the architecture choice
    depends on the bandwidth target.

  \item \textbf{Cascade beats native by an extra $Q_1$ if you can push
    $Q_1$.}  Comparing slopes within the cascade region itself,
    family B$_\text{I}$ (varying $Q_1$ at fixed $Q_2$) lies above
    B$_\text{II}$ (varying $Q_2$ at fixed $Q_1$) in the high-$Q_1$
    direction.  To make this quantitative, define the
    phase-bandwidth product $\Pi\equiv\phi_\text{NL}\cdot\Delta f$.
    At $Q_1=Q_2=Q$ (matched-$Q$ corner, $\gamma_1=\gamma_2$), using
    \eqref{eq:phi-A-scaling}--\eqref{eq:phi-B-scaling} and
    $\Delta f=\gamma_1/(2\pi)\propto 1/Q$:
    \begin{align*}
      \Pi^\text{A}\bigl|_{Q_1=Q}
          &\;\propto\;\chithree_\text{native}\,Q^{2}\cdot Q^{-1}
          \;=\;\chithree_\text{native}\,Q,\\
      \Pi^\text{B}\bigl|_{Q_1=Q_2=Q}
          &\;\propto\;\frac{|\chitwo|^{2}\,\eta_\text{geom}}{n_2^{\,2}}\,Q^{2}\,Q
              \cdot Q^{-1}
          \;=\;\frac{|\chitwo|^{2}\,\eta_\text{geom}}{n_2^{\,2}}\,Q^{2}.
    \end{align*}
    The cascade phase-bandwidth product at the matched-$Q$ corner
    scales as $Q^{2}\cdot R^{(0)}$, where
    $R^{(0)}\equiv|\chitwo|^{2}\eta_\text{geom}/(n_2^{\,2}\chithree_\text{native})$
    is the structural factor at unit $Q_2$.  The native architecture
    scales only as $Q$.  The cascade therefore beats native in
    phase-bandwidth by a factor of $Q$ at the matched corner -- one
    extra power of $Q$ relative to the slope-only comparison.
\end{enumerate}

% --------------------------------------------------------------
\subsection*{The crossover bandwidth}

Equating the native phase \eqref{eq:phi-A-scaling} to the cascade
phase \eqref{eq:phi-B-scaling} at the bandwidth-matched point
$\gamma_1=\gamma_2$ (the SH bottleneck balances the FM bandwidth)
and using the operating point $\Dom=2\gamma_2$ of
Section~\ref{sec:numerical} (so $|\chithree_\text{casc}|\sim
Q_2/(24\,n_2^{\,2})\cdot|\chitwo|^{2}\eta_\text{geom}$):
\begin{equation}
  \chithree_\text{native}\,Q_1^{2}\;\sim\;
  \frac{|\chitwo|^{2}\,\eta_\text{geom}}{24\,n_2^{\,2}}\,Q_1^{2}\,Q_2
  \;\Longrightarrow\;
  Q_2^\text{crossover}\;\sim\;
       \frac{24\,n_2^{\,2}\,\chithree_\text{native}}{|\chitwo|^{2}\,\eta_\text{geom}}.
  \label{eq:Q2-crossover}
\end{equation}
For LiNbO$_3$ at the values of Section~\ref{sec:numerical},
$Q_2^\text{crossover}\sim 3\!\times\!10^{2}$: any SH cavity with $Q_2$
above this threshold tips the comparison in favour of the cascade.
Since $Q_2 \sim 4000$ is routine in qBIC metasurfaces, LiNbO$_3$
\emph{always} prefers cascade in this regime.

For AlGaAs near the band edge, native $\chithree_\text{native}\sim
10^{-18}$~m$^{2}$/V$^{2}$ is enhanced by two-photon-resonance proximity
while $|\chitwo|\sim 10^{-10}$~m/V is only modestly larger than
LiNbO$_3$; \eqref{eq:Q2-crossover} gives $Q_2^\text{crossover}\sim
10^{5}$, which is above routinely achievable qBIC $Q$'s, so native
wins at all practical SH cavities.

For Si and Si$_3$N$_4$ the comparison is moot: $\chitwo\to 0$ by
inversion symmetry, and cascading is simply not available.  These
materials have to rely on native-Kerr resonance (or move to
periodically-poled cousin materials).

% --------------------------------------------------------------
\subsection*{Material survey}

Combining \eqref{eq:R-def} with the explicit form
\eqref{eq:chi3casc_explicit} at the nominal design point ($Q_2=
4{\times}10^{3}$, $\Dom=2\gamma_2$, $\eta_\text{geom}=0.3$, $n_2$ from
the material) gives the cascaded enhancement ratio for a few
representative materials:

\begin{center}
\begin{tabular}{lccccc}
\toprule
Material & $|\chitwo|$ [m/V] & $|\chithree_\text{native}|$ [m$^{2}$/V$^{2}$]
         & $n_2$ & $R$ at $Q_2{=}4{\times}10^{3}$ & $Q_2^\text{cross}$ \\
\midrule
LiNbO$_3$ (telecom)     & $5\!\times\!10^{-11}$   & $2\!\times\!10^{-21}$ & 2.26 & $\sim\!12$ & $\sim 3\!\times\!10^{2}$ \\
GaP (NIR)               & $1\!\times\!10^{-10}$   & $\sim\!10^{-19}$      & 3.1  & $\sim\!0.5$ & $\sim 8\!\times\!10^{3}$ \\
AlGaAs ($\sim$\,1550~nm)& $1\!\times\!10^{-10}$   & $\sim\!10^{-18}$      & 3.3  & $\sim\!0.05$ & $\sim 9\!\times\!10^{4}$ \\
AlN                     & $1\!\times\!10^{-11}$   & $\sim\!2\!\times\!10^{-21}$ & 2.1 & $\sim\!0.5$ & $\sim 7\!\times\!10^{3}$ \\
Si                      & $\equiv 0$              & $\sim\!4\!\times\!10^{-18}$ & 3.5  & \multicolumn{2}{c}{--- (no cascade)} \\
Si$_3$N$_4$             & $\sim 0$                & $\sim\!2.5\!\times\!10^{-21}$ & 2.0  & \multicolumn{2}{c}{--- (no cascade)} \\
\bottomrule
\end{tabular}
\end{center}

\noindent
Three regimes emerge.  In materials where the native $\chithree$ is
small relative to $|\chitwo|^{2}/(n_2^{2}\omega_1)$ -- chiefly
LiNbO$_3$ -- the cascade wins by an order of magnitude at routine
$Q_2$ and increasing $Q_2$ buys further linear gain.  In materials of
intermediate $\chithree_\text{native}$ such as GaP, the cascade and
native channels are comparable at $Q_2\sim 4\!\times\!10^{3}$, and the
architecture choice depends on whether one can push $Q_2$ above the
crossover \eqref{eq:Q2-crossover} into the cascade-dominated regime.
In materials where native $\chithree$ is enhanced by electronic
resonances near the band edge (AlGaAs, GaAs), the comparison flips:
native bulk-Kerr resonance is stronger at any practical $Q_2$, and
the engineering simplicity of a single resonance is a free bonus.
Si and Si$_3$N$_4$ have no cascade at all.  (At the
edge-of-adiabatic operating point $\Dom\!\sim\!\gamma_2$ of
Section~\ref{sec:numerical}, the $R$ values above double; the
qualitative picture is unchanged.)

% --------------------------------------------------------------
\subsection*{Caveats specific to this comparison}

\begin{enumerate}\setlength{\itemsep}{2pt}
  \item \emph{Cascade-specific loss.}  The cascaded TPA-like
        coefficient $\gamma_\text{eff}^{(2)}=|\beta|^{2}\gamma_2/\Dom^{2}$
        of Section~\ref{sec:adiab}(ii) has no analogue in low-loss native materials
        like LiNbO$_3$.  At fixed phase shift, cascade always carries
        more loss than native, by a factor $\gamma_2/(2\Dom)$ for the
        Kerr-to-loss ratio.  This penalises cascade in
        all-optical-switching contrast and in low-noise applications.

  \item \emph{Lower saturation.}  The cascade saturates at
        $|\beta||a_1|\sim|\Dom|$, beyond which the trivial fixed
        point becomes parametrically unstable
        (Section~\ref{sec:validity}, item~2).  At fixed
        $\Dom/\gamma_2$, the saturated FM amplitude
        $|a_1|_\text{sat}=|\Dom|/|\beta|\propto\gamma_2\propto 1/Q_2$
        (because $|\Dom|\propto\gamma_2\propto 1/Q_2$ while $|\beta|$
        is set by the modal overlap and is $Q_2$-independent), so
        the maximum useful stored energy
        $|a_1|^{2}_\text{sat}\propto 1/Q_2^{\,2}$.  Native-$\chithree$
        saturation, by contrast, is set by thermal/damage limits
        typically orders of magnitude higher.

  \item \emph{Fabrication complexity.}  Architecture A requires one
        high-$Q$ resonance at $\omega_1$; architecture B requires
        \emph{two} high-$Q$ resonances at $\omega_1$ and $2\omega_1$
        with simultaneously high spatial overlap.  Disorder broadens
        $\gamma_2$ inhomogeneously and can wash out the cascade-enhancement
        long before it affects a single-resonance native design.

  \item \emph{Speed.}  Native Kerr is essentially instantaneous
        ($\lesssim 10$~fs in dielectrics); cascade has bandwidth
        $\gamma_2$, i.e.\ ps-scale for routine $Q_2$.  For
        ultrafast applications (sub-ps), only native-Kerr is fast
        enough; cascade is locked out by its own sum rule.
\end{enumerate}

% --------------------------------------------------------------
\subsection*{Summary of the comparison}

\begin{center}
\begin{tabular}{p{0.27\textwidth}p{0.30\textwidth}p{0.36\textwidth}}
\toprule
Question & Answer & When it matters \\
\midrule
At fixed BW, who wins? &
Cascade by factor $R$; native if $R<1$. &
Sets the architecture decision for any narrowband application. \\

How does it scale with $Q$? &
$\phi^\text{A}\!\propto\!Q_1^{2}$; $\phi^\text{B}\!\propto\!Q_1^{2}Q_2$. &
Cascade has a free extra Q-knob; useful in high-Q regime. \\

What's the price of more $Q$? &
A: BW shrinks linearly; B: BW shrinks linearly in $Q_1$, also in $Q_2$. &
Both pay the sum rule; cascade more so. \\

Saturation? &
A: thermal/damage. B: $|\beta a_1|\sim|\Dom|$, much lower. &
Cascade is the limited one at high powers. \\

Ultrafast? &
A: $\sim$~fs; B: $\sim 1/\gamma_2$~ps. &
Use native for pulses $\lesssim 1$ ps. \\

Material choice? &
A: any high-$\chithree_\text{native}$ low-loss dielectric. B: needs $\chitwo\neq 0$. &
$\chitwo$-active materials only for cascade. \\
\bottomrule
\end{tabular}
\end{center}

\noindent
The take-away: the cascade is the right choice for \emph{narrowband,
CW-driven nonlinear photonics in $\chitwo$-active materials} (LiNbO$_3$,
GaP at telecom).  For Si/Si$_3$N$_4$ photonics, ultrafast switching,
high-power operation, or applications where two-photon loss is
intolerable, a native-$\chithree$ resonance is the right answer.

% =======================================================================
\section{Application landscape: Kerr-driven observables and devices}\label{sec:applications}

The previous section reduced the architecture question to a single
structural ratio
$R\equiv|\chithree_\text{casc}|/|\chithree_\text{native}|$.  But $R$
is only directly observable through a specific class of measurements:
those that depend on the FM-cavity Kerr coefficient
$\alpha_3^\text{eff}=\alpha_3+\alpha_3^\text{casc}$.  Different Kerr-driven
observables scale differently with $\alpha_3^\text{eff}$ -- linearly,
inversely, or quadratically -- and the cascade therefore appears in
each at a different "leverage".  In addition, several capabilities
exist that have \emph{no} native-Kerr analogue at all, and these are
where the cascaded route is uniquely useful rather than merely
quantitatively better.

We organise the discussion in three tiers.  Tier~1 collects
observables that admit a clean architecture-vs-architecture comparison
of the form (A+B)/(A)$=f(R)$.  Tier~2 collects capabilities specific
to the cascade with no native counterpart.  Tier~3 lists applications
that lie outside the single-cavity TCMT framework and require
extensions we do not develop here.

% --------------------------------------------------------------
\subsection{Tier 1: observables with quantitative \texorpdfstring{$R$}{R}-scaling}

% ............................................................
\paragraph{1. Self-phase modulation (SPM).}
Already established in Section~\ref{sec:comparison}.  On-resonance pump
at $\omega_p=\omega_1$ in the fully-radiative limit
$\gamma_1^{(\text{rad})}=\gamma_1$ gives a nonlinear phase
$\phi_\text{NL}=8\alpha_3^\text{eff}|s_+|^{2}/\gamma_1^{2}$ per input
photon \eqref{eq:phi-per-photon} (at the conventional critical-coupling
point $\gamma_1^{(\text{rad})}=\gamma_1/2$, $|a_1|^{2}=2|s_+|^{2}/\gamma_1$
and $\phi_\text{NL}=4\alpha_3^\text{eff}|s_+|^{2}/\gamma_1^{2}$; we
quote the lossless value to match \eqref{eq:phi-per-photon}).  Since $\phi_\text{NL}$ is linear
in $\alpha_3^\text{eff}=\alpha_3+\alpha_3^\text{casc}$, and using
the cascaded ratio \eqref{eq:R-def}, the architecture comparison gives
\begin{equation}
  \boxed{\;\frac{\phi_\text{NL}^{\text{A+B}}}{\phi_\text{NL}^{\text{A}}}
        =\frac{\alpha_3+\alpha_3^\text{casc}}{\alpha_3}
        =1+R\;.}
  \label{eq:SPM-ratio}
\end{equation}
Linear in $R$; the foundational observable.  Read off as
intensity-dependent transmission lineshape distortion or via
heterodyne phase measurement at the FM port.

% ............................................................
\paragraph{2. Cross-phase modulation (XPM).}
A separate weak signal beam, co-resonant with the FM cavity but
distinguishable by polarisation or by frequency offset
$|\Omega|<\gamma_1$, experiences a refractive-index shift driven by
the pump intensity $|a_1^\text{pump}|^{2}$.  In the cascade picture,
the pump generates an SH field that mediates DFG mixing onto the
signal; in the native picture, the bulk $\chithree$ provides the
mixing directly.  Both channels carry the same Boyd degeneracy factor
of $6$ (the index-permutation factor $2(2-\delta_{ij})$ for distinct
fields), so the relative architecture ratio inherits the SPM result:
\begin{equation}
  \boxed{\;\frac{\alpha_3^\text{XPM,\,A+B}}{\alpha_3^\text{XPM,\,A}}=1+R\;,\qquad
          \alpha_3^\text{XPM}=2\,\alpha_3^\text{SPM}\;.}
\end{equation}
The factor $2$ between XPM and SPM holds independently in both
architectures.  Experimentally: pump-probe transmission, with the
probe phase shift versus pump intensity read out via interferometry or
heterodyne.

\textbf{Important restriction.}  XPM in the strict sense requires a
second beam (distinguishable from the pump).  A single beam at
$\omega_p$ with a weak modulation at $\omega_p\pm\Omega$ instead drives
\emph{four-wave mixing} (item~4 below), not XPM, because the cascade
generates the conjugate sideband through DFG rather than a phase
shift.  The XPM coefficient as defined here applies to a genuinely
separate signal channel.

\textbf{Tensor-channel caveat.}  The $1+R$ ratio holds in the
\emph{scalar} TCMT picture used here.  In a real metasurface, the
cascade XPM goes through pump${\to}$SH${\to}$DFG-with-signal, which
involves $\chitwo_{ijk}$ contracted with mixed polarisation indices
($e{\cdot}e{\to}{\rm SH}$, then ${\rm SH}{\cdot}o{\to}{\rm signal}$ for
a $3m$ crystal with e-polarised pump and o-polarised signal, for
example).  These tensor contractions can differ from the
$\chitwo_{eee}$ that mediates SHG of the pump alone -- in some
polarisation channels they vanish by selection rule, in others they
are comparable.  The native-XPM channel by contrast uses
$\chithree_{eooe}$ or similar, which is governed by the host
$\chithree$ tensor.  The $1+R$ ratio holds for tensor channels in
which both cascade and native pieces are non-vanishing and
geometrically similar; in mixed-channel experiments the relative
factor may differ.  A practical experiment should verify the relevant
tensor contraction for the intended polarisation geometry.

% ............................................................
\paragraph{3. Bistability threshold and all-optical switching energy.}
The steady-state amplitude equation at pump detuning
$\delta_p\equiv\omega_p-\omega_1$, obtained by setting $\dot a_1=0$ in
\eqref{eq:a1} (with $a_2$ adiabatically eliminated and the cascaded
TPA neglected for the cusp analysis), is
\begin{equation}
  \bigl[\tfrac{\gamma_1}{2}-i\bigl(\delta_p+\alpha_3^\text{eff}|a_1|^{2}\bigr)\bigr]\,a_1
   \;=\;\sqrt{\gamma_1^\text{rad}}\,s_+.
  \label{eq:bist-SS}
\end{equation}
Taking moduli squared and writing $y\equiv\alpha_3^\text{eff}|a_1|^{2}/(\gamma_1/2)$,
$x\equiv\delta_p/(\gamma_1/2)$, $u\equiv|\alpha_3^\text{eff}|\,|s_+|^{2}\,
\gamma_1^\text{rad}/(\gamma_1/2)^{3}$ (a dimensionless input intensity),
\eqref{eq:bist-SS} becomes the standard Kerr-bistability cubic
$y[1+(x-y)^{2}]=u$.  This cubic has three real roots
$y_-<y_\text{mid}<y_+$ -- the bistable regime -- iff the discriminant
condition is satisfied, which is well known
\cite{DrummondWalls1980,Boyd} to require
\begin{equation}
  \mathrm{sign}(\delta_p)=-\mathrm{sign}(\alpha_3^\text{eff}),
  \qquad
  \frac{|\delta_p|}{\gamma_1/2}>\sqrt{3}\;\Longleftrightarrow\;
  |\delta_p|>\frac{\sqrt{3}}{2}\,\gamma_1.
  \label{eq:bist-cusp}
\end{equation}
The sign condition is automatic: a positive (self-focusing)
$\alpha_3^\text{eff}>0$ pulls the cavity resonance down with growing
intensity, so a red-detuned pump ($\delta_p<0$) gets pulled into
resonance and bistability ensues; for a self-defocusing
$\alpha_3^\text{eff}<0$ (e.g.\ a cascade with $\Dom>0$), the same
analysis with the opposite sign of $\delta_p$ applies.  The
$|\delta_p|>\sqrt{3}\gamma_1/2$ inequality is the cusp condition.

The snap-up power follows from the saddle-node bifurcation of the
cubic.  At the cusp $|\delta_p|=\sqrt{3}\gamma_1/2$, the critical
intracavity intensity is $|a_1|^{2}_\text{c}=|\delta_p|/|\alpha_3^\text{eff}|
=\sqrt{3}\gamma_1/(2|\alpha_3^\text{eff}|)$, and the corresponding
input power is
\begin{equation}
  |s_+|^{2}_\text{bist}\;\propto\;
     \frac{\gamma_1^{2}\,|\delta_p|}{|\alpha_3^\text{eff}|\,\gamma_1^\text{rad}}
     \;\propto\;\frac{\gamma_1^{2}}{|\alpha_3^\text{eff}|}
     \quad\text{(at fixed }\delta_p/\gamma_1\text{).}
  \label{eq:bist-threshold}
\end{equation}
Inversely proportional to the Kerr coefficient, so the architecture
ratio is inverted relative to the phase shift:
\begin{equation}
  \boxed{\;\frac{|s_+|^{2,\,\text{A+B}}_\text{bist}}{|s_+|^{2,\,\text{A}}_\text{bist}}=\frac{1}{1+R}\;.}
  \label{eq:bist-ratio}
\end{equation}
For LiNbO$_3$ ($R\sim 12$ at the nominal $\Dom=2\gamma_2$,
Section~\ref{sec:numerical}), the cascade lowers the bistability
threshold by $\sim\!13\times$ (and by $\sim\!26\times$ at the
edge-of-adiabatic $\Dom\sim\gamma_2$).  Re-framed as an all-optical
switching energy: biasing the system on the lower branch and applying
a control pulse to push past the saddle-node, the minimum required
control energy obeys
\begin{equation}
  \boxed{\;\frac{E_\text{switch}^{\text{A+B}}}{E_\text{switch}^{\text{A}}}=\frac{1}{1+R}\;.}
\end{equation}
The switching speed is set by the FM ringdown $1/\gamma_1$ (ps for
$Q_1=10^{3}$); the readout is a direct power measurement on a
photodiode.  \textbf{This is the canonical photonic-computing
application.}

% ............................................................
\paragraph{4. Four-wave-mixing (FWM) wavelength conversion.}
Driving with a strong pump at $\omega_p\approx\omega_1$ and a weak
signal at $\omega_s=\omega_p-\Omega$ generates an idler at
$\omega_i=2\omega_p-\omega_s=\omega_p+\Omega$.  In the cascade
mechanism: the pump generates SH at $2\omega_p\approx\omega_2$, which
then DFG-mixes with the signal to produce the idler at the FM cavity.

\paragraph*{Derivation of the conversion efficiency.}
Write the FM amplitude as $a_1=a_p+a_s\,e^{-i\Omega t}+a_i\,e^{+i\Omega t}$
in the frame rotating at $\omega_p$, with $|a_s|,|a_i|\ll|a_p|$.
Substituting into the effective FM equation
$\dot a_1 = (-\gamma_1/2)\,a_1+i\,\alpha_3^\text{eff}|a_1|^{2}a_1
+\sqrt{\gamma_1^\text{rad}}s_+$ and collecting the $e^{+i\Omega t}$
Fourier component generates an idler source
\begin{equation*}
  \mathcal{S}_i\;=\;i\,\alpha_3^\text{eff}\,a_p^{2}\,a_s^{*},
\end{equation*}
which is the standard FWM source proportional to pump${}^{2}\times$
signal-conjugate.  Expanding $|a_1|^{2}a_1$ to first order in the
small fields and collecting the $e^{+i\Omega t}$ Fourier component
gives two pieces: $a_p^{2}\,a_s^{*}$ (the FWM source above) and a
$2|a_p|^{2}a_i$ piece that renormalises the idler detuning -- a
pump-induced XPM self-shift that adds a real part
$\alpha_3^\text{eff}\,2|a_p|^{2}$ to the idler resonance condition.
In the small-signal estimate below we drop this self-shift (it is
$O(\alpha_3^\text{eff}|a_p|^2/\gamma_1)$ smaller than the propagator
denominator $\gamma_1/2$ once the pump is below the bistability
threshold of item~3).  The idler steady-state response is then
\begin{equation*}
  a_i \;=\; \frac{i\,\alpha_3^\text{eff}\,a_p^{2}\,a_s^{*}}{\gamma_1/2-i\Omega}
       \;\stackrel{|\Omega|<\gamma_1/2}{\simeq}\;
       \frac{i\,\alpha_3^\text{eff}\,a_p^{2}\,a_s^{*}}{\gamma_1/2},
\end{equation*}
and the idler output via the input--output relation
$s_-^{(i)}=\sqrt{\gamma_1^\text{rad}}\,a_i$ gives an idler power
$|s_-^{(i)}|^{2}=\gamma_1^\text{rad}|a_i|^{2}$ while the signal input
generates an intracavity signal
$|a_s|=2\sqrt{\gamma_1^\text{rad}}|s_+^{(s)}|/\gamma_1$.  Combining,
\begin{equation}
  \eta_\text{FWM}\;\equiv\;\frac{|s_-^{(i)}|^{2}}{|s_+^{(s)}|^{2}}
  \;\propto\;\frac{|\alpha_3^\text{eff}|^{2}\,|a_p|^{4}}{(\gamma_1/2)^{2}}\,,
  \label{eq:eta-FWM}
\end{equation}
\emph{quadratic} in the effective Kerr coefficient.  The architecture
ratio is therefore
\begin{equation}
  \boxed{\;\frac{\eta_\text{FWM}^{\text{A+B}}}{\eta_\text{FWM}^{\text{A}}}=(1+R)^{2}\;.}
\end{equation}
For LiNbO$_3$ ($R\sim 12$ at the nominal $\Dom=2\gamma_2$), this is a
factor of $\sim 170$ in conversion efficiency at the conservative
operating point, and approaches $\sim 700$ if one operates at
$\Dom\sim\gamma_2$.  Either way, FWM is the most strongly
differentiating single-observable test of the cascade mechanism.

\textbf{Bandwidth constraint.}  All three frequencies must lie within
$\min(\gamma_1,\gamma_2)$.  For $Q_1\ll Q_2$ the FM linewidth bounds
$\Omega$ via the $1/(\gamma_1/2-i\Omega)$ factor in $a_i$; for
$Q_1\gtrsim Q_2$ the cascade bandwidth $\gamma_2$ takes over via the
AC response \eqref{eq:alpha3-AC} of $\alpha_3^\text{eff}(\Omega)$.
Cascade-FWM is therefore restricted to a frequency window
$\Omega\lesssim\min(\gamma_1,\gamma_2)/(2\pi)\sim$~GHz at telecom
$Q$'s.

% --------------------------------------------------------------
\subsection{Tier 2: cascade-only capabilities}

These applications exploit aspects of the cascaded $\chithree_\text{eff}$
that have no analogue in a native-$\chithree$ resonance: the
$\Dom$-dependence of the coefficient itself, and the imaginary part
$\gamma_3^\text{eff}=|\beta|^{2}(\gamma_2/2)/[\Dom_\text{SH}^{2}+(\gamma_2/2)^{2}]$
which produces an intensity-dependent loss whose magnitude is
\emph{also} $\Dom$-tunable.  Architecture A simply does not provide
either knob.

% ............................................................
\paragraph{5. Sign-tunable saturable absorber.}
Throughout this item we write $\Dom_\text{SH}\equiv\Dom=2\omega_1-\omega_2$
for the SH detuning, to distinguish it visually from the FM pump
detuning $\delta_p$ that appears in items 3 and 4 above; numerically
they are the same quantity used since \S\ref{sec:tcmt}.

Operating near the SH resonance ($|\Dom_\text{SH}|\lesssim\gamma_2$),
take the imaginary part of \eqref{eq:chi3casc-tensor}.  Writing
$1/(\Dom+i\gamma_2/2)=(\Dom-i\gamma_2/2)/[\Dom^{2}+(\gamma_2/2)^{2}]$,
the imaginary part of $\chithree_\text{casc}$ is
\begin{equation}
  \mathrm{Im}\,\chithree_\text{casc}(\Dom)
   \;=\;\frac{\eps_0\,\omega_1}{6}\cdot
        \frac{|M|^{2}\,(\gamma_2/2)}{[\Dom^{2}+(\gamma_2/2)^{2}]\,\int|\f_1|^{4}\dthreer}.
  \label{eq:Im-chi3casc}
\end{equation}
The corresponding TCMT loss coefficient (obtainable directly by
keeping the imaginary part of the adiabatic-elimination step in
Section~\ref{sec:adiab} without expanding for small $\gamma_2/\Dom$)
is
\begin{equation}
  \gamma_3^\text{eff}(\Dom_\text{SH})
   \;=\;\frac{|\beta|^{2}\,(\gamma_2/2)}{\Dom_\text{SH}^{2}+(\gamma_2/2)^{2}},
  \label{eq:gamma3-Lorentzian}
\end{equation}
which is precisely $\gamma_\text{eff}^{(2)}=|\beta|^{2}\gamma_2/\Dom^{2}$
of Section~\ref{sec:adiab}(ii) in the adiabatic limit $|\Dom|\gg\gamma_2$,
and rises to a Lorentzian peak at $\Dom_\text{SH}=0$:
\begin{equation}
  \gamma_3^\text{eff}(\Dom_\text{SH}=0)\;=\;\frac{|\beta|^{2}}{\gamma_2/2}
  \;=\;\frac{2\,Q_2}{\omega_2}\,|\beta|^{2}
  \qquad(\text{Lorentzian, FWHM}=\gamma_2).
  \label{eq:gamma3-peak}
\end{equation}
In the language of an effective material parameter, using
\eqref{eq:Im-chi3casc} and the same $1/n_2^{\,2}$ analysis as in
Section~\ref{sec:adiab} (replacing $\Dom$ by $\gamma_2/2$ in the
denominator of the explicit form \eqref{eq:chi3casc_explicit}),
\begin{equation}
  \mathrm{Im}\,\chithree_\text{casc}\big|_{\Dom_\text{SH}=0}
       \;\sim\;\frac{\omega_1}{3\,n_2^{\,2}\,\gamma_2}\,|\chitwo|^{2}\,\eta_\text{geom}
       \;=\;\frac{Q_2}{6\,n_2^{\,2}}\,|\chitwo|^{2}\,\eta_\text{geom},
  \label{eq:Im-chi3-explicit}
\end{equation}
where we used $\omega_1/\gamma_2 = Q_2/2$.  Compared with the real
part at $\Dom=\gamma_2/2$ (where $|\mathrm{Re}\,\chithree_\text{casc}|
=Q_2/(12 n_2^{\,2})\,|\chitwo|^{2}\eta_\text{geom}$ from
\eqref{eq:chi3casc_explicit}), the peak imaginary part is a factor of
$2$ larger -- the same Lorentzian peak-to-shoulder ratio that maps the
dispersive real part to the absorptive imaginary part.
This is a \emph{saturable absorption} channel (intensity-dependent
loss) that, unlike intrinsic two-photon absorption in semiconductors,
is \emph{detunable} away from resonance: shifting $\omega_2$ by
$\gtrsim$ a few $\gamma_2$ (electro-optically, thermo-optically, or
by strain) suppresses the loss by $1/\Dom_\text{SH}^{2}$ from
\eqref{eq:gamma3-Lorentzian}.

Comparison with SESAM (semiconductor saturable absorber mirror, the
workhorse of pulsed-laser mode locking \cite{Keller1996}): SESAMs have
a fixed, material-determined recovery time set by excited-state
population dynamics; cascaded SA has recovery time $\sim 1/\gamma_2$
(ps for
$Q_2=4{\times}10^{3}$) and is all-optical, with no excited-state
physics or thermal lifetime drift.  The trade-off is that cascaded SA
requires a $\chitwo$-active substrate and double-resonance fabrication.

% ............................................................
\paragraph{6. All-optical amplitude modulator via $\Dom$ tuning.}
Modulating $\Dom(t)=\Dom_0+\delta\Omega\cos(\Omega_\text{mod}t)$ via
electro-optic, thermo-optic, or piezo-electric actuation of the SH
cavity periodically modulates $\alpha_3^\text{eff}$, and therefore
the FM transmission $|s_-/s_+|^{2}$ at fixed pump.  For small
modulation $\delta\Omega\ll|\Dom_0|$ and $|\Dom_0|\gg\gamma_2/2$,
the adiabatic-limit derivative
\begin{equation}
  \frac{\partial\alpha_3^\text{casc}}{\partial\Dom}\bigg|_{\Dom_0}
   \;\simeq\;\frac{|\beta|^{2}}{\Dom_0^{2}}
  \label{eq:dalpha-dDom}
\end{equation}
gives, in the DC limit ($\Omega_\text{mod}\to 0$), a fractional
modulation $\delta\alpha_3^\text{casc}/\alpha_3^\text{casc}=\delta\Omega/\Dom_0$:
a fractional SH-cavity detuning produces an equal fractional change
in $\alpha_3^\text{casc}$.

\paragraph*{Transfer function: low-pass plus AC band-pass.}
For finite $\Omega_\text{mod}$ the cascaded coefficient cannot track
the SH-cavity detuning instantaneously.  The correct response must be
obtained from the \emph{time-dependent} SH equation with modulated
detuning -- linearising
$\Dom(t)=\Dom_0+\delta\Omega\cos(\Omega_\text{mod}t)$ in $\delta\Omega$
and solving the SH equation in steady state gives, for the
\emph{positive-frequency sideband} of $\delta a_2$,
\begin{equation}
  A_+(\Omega_\text{mod})
   \;=\;\frac{(\delta\Omega/2)\,a_{2,0}}{\gamma_2/2+i(\Omega_\text{mod}-\Dom_0)},
  \qquad a_{2,0}=\frac{-\beta\,a_1^{2}}{\Dom_0+i\gamma_2/2},
  \label{eq:H-modulator}
\end{equation}
which is resonant at $\Omega_\text{mod}=\Dom_0$ (the sideband of
$\Dom(t)$ matches the SH-cavity resonance in the rotating frame).
The fractional $a_2$ modulation, which carries directly into the
cascade contribution to $\alpha_3^\text{eff}$, is
\begin{equation}
  \left|\frac{\delta a_2}{a_{2,0}}\right|_{\Omega_\text{mod}}
   \;=\;\frac{\delta\Omega/2}{\sqrt{(\gamma_2/2)^{2}+(\Omega_\text{mod}-\Dom_0)^{2}}}.
\end{equation}
Its values at the two limits are
\begin{equation}
  \left|\frac{\delta a_2}{a_{2,0}}\right|_\text{peak\,at\,}\Omega=\Dom_0
   \;=\;\frac{\delta\Omega}{\gamma_2},
  \qquad
  \left|\frac{\delta a_2}{a_{2,0}}\right|_\text{DC}
   \;=\;\frac{\delta\Omega}{\sqrt{\Dom_0^{2}+(\gamma_2/2)^{2}}},
\end{equation}
so the AC-peak enhancement over DC is
\begin{equation}
  \boxed{\;
  \frac{|\delta\alpha_3^\text{eff}|_\text{peak\,at\,}\Omega_\text{mod}=\Dom_0}
       {|\delta\alpha_3^\text{eff}|_\text{DC}}
  \;=\;\frac{\sqrt{\Dom_0^{2}+(\gamma_2/2)^{2}}}{\gamma_2}
  \;\xrightarrow[\Dom_0\gg\gamma_2]{}\;\frac{|\Dom_0|}{\gamma_2}.\;}
  \label{eq:AC-peak-enhancement}
\end{equation}
At the conservative operating point $\Dom_0=2\gamma_2$ the AC peak is
$\sqrt{4.25}\approx 2.06\times$ the DC modulation depth -- a modest
but useful boost.  This number agrees with the TCMT measurement
($1.92\times$ in the AM-modulator figure of the simulation
report).  Note that this is the boost on the cascade-side
$\alpha_3^\text{eff}$ modulation amplitude; the FM-port transmission
modulation also picks up a factor $1/|\gamma_1/2 + i\Omega_\text{mod}|$
from the FM-cavity filtering, which is close to unity at
$\Omega_\text{mod}\!=\!2\gamma_2$ when $\gamma_1\gg 2\gamma_2$ (the
worked-example regime).

The transfer function from $\delta\Omega(t)$ to the FM transmission
modulation therefore has two features:
\begin{itemize}\setlength{\itemsep}{1pt}
\item a low-pass response with roll-off at $\Omega_\text{mod}\sim\gamma_2/2$
      (single-pole, set by the cascade bandwidth);
\item a Lorentzian band-pass peak at $\Omega_\text{mod}=\Dom_0$ of
      FWHM $\gamma_2$ and peak-to-DC ratio $4\Dom_0^{2}/\gamma_2^{2}$.
\end{itemize}
At telecom $Q_2=4{\times}10^{3}$ this gives a DC bandwidth
$\lesssim\gamma_2/(2\pi)\!\simeq\!97$~GHz and an AC operating point at
$\Omega_\text{mod}/(2\pi)=|\Dom_0|/(2\pi)\simeq 200$~GHz when
$\Dom_0=2\gamma_2$.

This is a genuinely novel capability: native-Kerr cavities have no
mechanism for converting a detuning modulation into an amplitude
modulation, since the bulk $\chithree$ is not detuning-dependent.

% ............................................................
\paragraph{7. Passive optical limiter via the parametric ceiling.}
The cascade saturation $|\beta||a_1|\to|\Dom|$
(Section~\ref{sec:validity}, item~2) is more than a small-signal-breakdown
caveat: it is a \emph{hard, passive intensity clamp} at intracavity
energy $|a_1|^{2}_\text{cap}\simeq(\Dom/\beta)^{2}$.  Above this
threshold the trivial fixed point destabilises and the FM clamps at
the cap regardless of further pump increase, with the excess input
dissipated through the cascaded TPA channel.  In an integrated
metasurface this is an all-optical \emph{power limiter} -- the same
function that absorption-based limiters serve in laser safety glasses,
but realised dispersively with no thermal load on the absorbing layer
(loss is in the SH cavity, which can be heat-sunk separately).
The ceiling is tunable by adjusting $\Dom$: at the design point
$\Dom=2\gamma_2$ the clamp is set; sweeping $\Dom$ tunes the ceiling
quadratically.  Native-$\chithree$ resonance has no such ceiling --
intensity scales with input up to material damage.  See
\S\ref{sec:adiab}--\S\ref{sec:validity} for the underlying mechanism.

% ............................................................
\paragraph{8. $\Dom$-readout transducer for sensing.}
Inverting the previous: any process that shifts $\omega_2$
(temperature, strain, applied electric field via Pockels effect on
the SH cavity, adsorbed molecules, etc.)\ shifts $\Dom$, which shifts
$\alpha_3^\text{eff}$, which shifts the SPM phase $\phi_\text{NL}$.
The cascade thus provides a built-in transduction from arbitrary
$\omega_2$-perturbations to a Kerr observable.

The sensitivity follows from the chain rule
$d\phi_\text{NL}/d\Dom=(8|s_+|^{2}/\gamma_1^{2})\cdot(d\alpha_3^\text{casc}/d\Dom)$
(using $\phi_\text{NL}\propto\alpha_3^\text{eff}$ from
\eqref{eq:phi-per-photon}).  Differentiating the Lorentzian
$\alpha_3^\text{casc}(\Dom)=-|\beta|^{2}\Dom/[\Dom^{2}+(\gamma_2/2)^{2}]$
of Section~\ref{sec:adiab} explicitly (with $b\equiv\gamma_2/2$):
\begin{equation}
  \frac{d\alpha_3^\text{casc}}{d\Dom}\;=\;-|\beta|^{2}\,
        \frac{b^{2}-\Dom^{2}}{[\Dom^{2}+b^{2}]^{2}}\,.
  \label{eq:dadDom}
\end{equation}
Solving $d|d\alpha/d\Dom|/d\Dom=0$ gives \emph{two} extrema: a global
maximum at $\Dom=0$ and a smaller local maximum at $|\Dom|=b\sqrt{3}$
(factor of $8$ below the global one).  The peak slope is therefore
\begin{equation}
  \boxed{\;
    \left|\frac{d\alpha_3^\text{casc}}{d\Dom}\right|_{\Dom=0}
     \;=\;\frac{|\beta|^{2}}{b^{2}}
     \;=\;\frac{4\,|\beta|^{2}}{\gamma_2^{\,2}}
     \;\propto\;Q_2^{\,2}\,,
  \;}
  \label{eq:dadDom-max}
\end{equation}
giving
\begin{equation}
   \left|\frac{d\phi_\text{NL}}{d\Dom}\right|_\text{max}
     \;=\;\frac{32\,|s_+|^{2}\,|\beta|^{2}}{\gamma_1^{\,2}\,\gamma_2^{\,2}}
     \;\propto\;Q_1^{\,2}\,Q_2^{\,2}\,.
\end{equation}
\textbf{The sensor sensitivity scales as} $Q_2^{\,2}$, not $Q_2$ --
quadratically more sensitive in the SH cavity factor than the SPM
phase shift itself, and gated by the FM cavity $Q_1^{\,2}$ in the
usual way.  The extra power of $Q_2$ comes from the Lorentzian slope:
a coefficient of amplitude $\sim |\beta|^{2}/\gamma_2$ varies over a
width $\sim\gamma_2$, giving a derivative $\sim|\beta|^{2}/\gamma_2^{\,2}$.

\textbf{Operating-point caveat.}  $\Dom=0$ is also exactly where the
cascaded two-photon-like loss
$\gamma_3^\text{eff}=|\beta|^{2}b/[\Dom^{2}+b^{2}]$ peaks
(eq.~\eqref{eq:gamma3-Lorentzian}) -- the sensor's most sensitive
point is its lossiest point.  Useful figures of merit therefore
depend on the weighting: the ratio (slope)$/\gamma_3^\text{eff}$ is
monotonic in $|\Dom|$ within each Lorentzian branch and asymptotes
to $2/\gamma_2$ at both $\Dom\to 0$ and $|\Dom|\to\infty$; it has a
minimum (zero) at $|\Dom|=b$ where the slope itself vanishes.  In
practice the sensor should be biased a few $\gamma_2$ off-resonance
-- enough to suppress TPA by $\sim(\gamma_2/\Dom)^{2}$ while
accepting a slope smaller than \eqref{eq:dadDom-max} by a similar
factor.  The fundamental $Q_2^{\,2}$ scaling survives this
trade-off; the prefactor is reduced by an $O(1)$ amount depending
on the chosen $\Dom/\gamma_2$.

\textbf{Shot-noise-limited figure of merit.}  The $Q_2^{2}$ scaling
applies to the raw slope $d\alpha_3^\text{casc}/d\Dom$.  For a
shot-noise-limited readout the relevant figure of merit is
slope$/\sqrt{\gamma_3^\text{eff}}$ (the loss sets the noise floor
through the detected-photon shot noise), which scales as
$Q_2/\sqrt{Q_2}=Q_2^{3/2}$.  Either scaling is faster than the linear
$Q_2$ enhancement of the cascaded Kerr coefficient itself.  Pushing
$Q_2$ continues to pay off, but the practical FoM exponent depends on
which noise channel dominates.

For a representative SH-cavity transducer in LiNbO$_3$
($dn_e/dT\!\simeq\!4\!\times\!10^{-5}$/K~\cite{Nikogosyan2005} gives a
bare $d\omega_2/dT\!\simeq\!-\omega_2\,(dn/dT)/n\!\sim\!2\pi\times
7$~GHz/K; the meta-atom-mode coupling reduces this by the mode-volume
overlap with the LiNbO$_3$, typically a factor $\!\sim\!3$--$5$, giving
an effective $\!\sim\!2\pi\times 1$--$2$~GHz/K), with
$\gamma_2/2\pi\!\sim\!100$~GHz at $Q_2=4{\times}10^{3}$, even after the
TPA-penalised bias this resolves mK-scale temperature shifts at modest
pump power and standard photodetector noise.  Applications
include refractometric biosensing, strain mapping, integrated
thermometry, and (with appropriate cavity design) electric-field
sensing via the Pockels effect.

\textbf{This may be the most useful application of the cascade
framework}, because it converts a familiar Kerr-readout chain into a
sensor whose figure of merit scales with $Q_2^{\,2}$ rather than the
single-power-of-$Q$ ceiling familiar from linear-cavity sensors.

% --------------------------------------------------------------
\subsection{Tier 3: applications outside the present scope}

The following applications use the cascade but require theoretical
extensions we do not develop:

\begin{itemize}\setlength{\itemsep}{2pt}
  \item \emph{Pulse compression and GVD compensation.}  The
    sign-tunable cascaded Kerr can compensate either sign of
    group-velocity dispersion in subsequent free-space or fibre
    propagation, leading to soliton-like compression or anti-soliton
    chirp.  This is a \emph{propagation-level} effect; the metasurface
    provides the local nonlinearity, but characterisation requires
    convolution with downstream linear optics.  See \cite{Liu1999,Bache2007}.

  \item \emph{Optical memory / flip-flop.}  Same physics as
    all-optical switching (item~3 above); the "memory time" is the
    upper-branch dwell time, limited by thermo-optic drift,
    fabrication detuning, and slow background fluctuations rather than
    by anything in the TCMT model.  Not a single-cavity TCMT effect to
    formalise.

  \item \emph{Logic gates.}  Two co-resonant control inputs driving
    one bistable element implement AND/OR by threshold combinations.
    This is a multi-port extension of bistability and requires the
    multi-port input-output theory of \cite{SuhWangFan2004} rather
    than the single-port TCMT used here.

  \item \emph{Degenerate optical parametric oscillation (OPO) and
    squeezed-light generation.}  Above the parametric threshold the
    cascade system \emph{is} an OPO, with above-threshold conversion
    efficiency and below-threshold squeezing of the appropriate
    quadrature.  But OPO operation requires substantial
    $\gamma_2^{(\text{rad})}$ at the SH port (to extract the
    parametric output), which contradicts the cascade design choice
    of suppressing SH out-coupling.  These are therefore
    \emph{different devices} with conflicting design constraints
    rather than alternative uses of the same metasurface.  Squeezed
    light additionally requires quantum input-output theory beyond the
    classical TCMT framework adopted in this note.

  \item \emph{Kerr frequency combs.}  Single-mode cascaded $\chithree$
    can in principle support a Lugiato--Lefever-type comb, but this
    requires a multi-mode generalisation of the TCMT (one resonance
    per comb tooth) together with dispersion engineering of the FM
    family.  A substantial separate project.

  \item \emph{Phase conjugation.}  The DFG step in the cascade is by
    construction phase-conjugate to the FM input.  This is implicit in
    the dynamics of $a_2$ and the back-action onto $a_1$; making it
    explicit requires tracking the phase of the SH amplitude through
    the cascade and is best handled in a dedicated treatment.
\end{itemize}

% --------------------------------------------------------------
\subsection{Summary of observables}

\begin{center}
\renewcommand{\arraystretch}{1.2}
\begin{tabular}{lcl}
\toprule
Observable & (A+B)/(A) ratio & Comment \\
\midrule
SPM phase shift                  & $1+R$       & Foundational \\
XPM phase shift                  & $1+R$       & Same $R$; requires separate signal \\
Bistability threshold $|s_+|^{2}$ & $1/(1+R)$  & Inverted scaling \\
All-optical switching energy     & $1/(1+R)$   & Photonic-computing application \\
FWM conversion efficiency        & $(1+R)^{2}$ & Most differentiating \\
\midrule
Saturable absorber rate          & n/a         & Cascade-only; $\Dom$-tunable \\
AM modulator response            & n/a         & Cascade-only; $\Dom(t)$-driven \\
Optical-limiter clamp $|a_1|^{2}_\text{cap}$ & $\sim(\Dom/\beta)^{2}$ &
   Cascade-only; passive intensity ceiling \\
Sensor sensitivity               & $\propto Q_2^{\,2}$ & Cascade-only; $\Dom$-readout, biased a few $\gamma_2$ off-resonance \\
\bottomrule
\end{tabular}
\end{center}

\noindent
Tier 1 establishes the cascade as a quantitative replacement for
bulk-Kerr-driven nonlinear optics, with leverage factors of
$1+R$ to $(1+R)^{2}$ depending on observable.  Tier 2 establishes the
cascade as a source of \emph{new} capabilities -- saturable
absorption, amplitude modulation, intensity limiting, and sensing --
that the native-$\chithree$ architecture cannot provide at any $Q_1$.
Together the two tiers cover most of the experimental opportunities a
doubly-resonant cascaded $\chithree$ metasurface opens up.

% =======================================================================
\section{Validity, caveats, and what to verify}\label{sec:validity}

\begin{enumerate}
  \item \textbf{Adiabatic-elimination breakdown.} The naive form
        \eqref{eq:chi3casc} ($\chithree_\text{casc}\propto 1/\Dom$)
        diverges as $\Dom\to 0$; the full Lorentzian
        $-|\beta|^2/(\Dom+i\gamma_2/2)$ retained in
        \eqref{eq:chi3casc-tensor} does not.  The cascaded-Kerr
        description (\emph{adiabatic} elimination of $a_2$) is strictly
        valid for $|\Dom|\gtrsim 2\gamma_2$; we quote results down to
        $|\Dom|\sim\gamma_2$ as an extrapolation that doubles the
        enhancement (real-part Lorentzian peaks at $|\Dom|=\gamma_2/2$,
        so $R$ is bounded above by twice its $\Dom=2\gamma_2$ value).
        At $|\Dom|=\gamma_2$ the loss-to-Kerr ratio
        $\gamma_3^\text{eff}/|\alpha_3^\text{eff,casc}|=\gamma_2/(2|\Dom|)$
        reaches $1/2$ and the slow-FM-modulation
        ($|\dot a_1|\ll|\Dom|\,a_1$) requirement also tightens; below
        $|\Dom|\!\sim\!\gamma_2$ one must keep $a_2$ as a dynamical
        variable.  On resonance ($\Dom=0$) the steady-state SH amplitude
        is bounded by $|a_2|=|\beta||a_1|^{2}/(\gamma_2/2)$.

  \item \textbf{Pump depletion / saturation.} The derivation is
        small-signal: it assumes $|a_2|\ll|a_1|$, equivalently
        $|\beta||a_1|^{2}/|\Dom|\ll|a_1|$, i.e.\
        $|\beta||a_1|\ll|\Dom|$.  Above this threshold the SH is no
        longer slaved to the FM and the response saturates; in fact,
        exactly at the saturation point one finds an analogue of the
        parametric oscillation threshold.  The maximum useful FM
        amplitude is therefore $|a_1|_\text{sat}=|\Dom|/|\beta|$,
        which scales as $1/Q_2$ at fixed $\Dom/\gamma_2$ (because
        $|\Dom|\propto\gamma_2\propto 1/Q_2$ while $|\beta|$ is set
        by the modal overlap, independent of $Q_2$); equivalently,
        the maximum stored energy $|a_1|^{2}_\text{sat}\propto
        1/Q_2^{\,2}$.  High-$Q_2$ devices saturate at lower
        powers.

  \item \textbf{Two-photon-like loss.} The imaginary part of
        $\alpha_3^\text{eff}$ -- the cascaded two-photon-absorption
        analogue -- is unavoidable and bounded below by the same
        overlap.  It does \emph{not} dump energy outside the system in
        the limit $\gamma_2^\text{rad}\to 0$, but rather absorbs into
        the SH mode and from there into the material.  This is
        identical to two-photon absorption from an experimental
        standpoint and limits the Kerr-driven phase shift before
        bleaching.

  \item \textbf{Thermo-optic and free-carrier nonlinearities} in
        dielectric metasurfaces are notoriously easy to confuse with
        ``Kerr''.  At the field intensities required to see a
        measurable cascaded $\chithree$ phase shift (several
        MW/cm$^{2}$ in the resonator), thermal effects in absorbing
        materials (Si, GaAs near the band edge) often dominate.
        LiNbO$_3$ is comparatively benign in this respect.

  \item \textbf{Tensor-overlap evaluation.} The dimensionless overlap
        $\eta_\text{geom}$ is assumption-driven in this note.  It must
        be evaluated by full-wave simulation (e.g.\ COMSOL or MEEP
        eigenmode solver, post-processing the tensor contraction in
        the volume of the meta-atoms).  The $\sim 0.3$ figure is
        plausible for a properly engineered qBIC pair but is by no
        means generic -- random mode pairings typically give
        $\eta_\text{geom}$ of order $10^{-3}$ or less, especially when
        the SH-mode parity does not match the symmetric tensor
        contraction $\f_1\f_1$ at unit-cell level.

  \item \textbf{Avoid the ``single $n$'' simplification.} Use
        $\epstilde(\ri,\omega_1)$ and $\epstilde(\ri,2\omega_1)$
        separately in the energy normalisation \eqref{eq:Brillouin}.
        The dispersion of $\chitwo$ between $\omega_1$ and $2\omega_1$
        is a few percent in LiNbO$_3$ but should be carried explicitly
        when comparing with experiment.

  \item \textbf{Disorder and inhomogeneous broadening.}  The $Q_2$
        quoted throughout is the \emph{element} quality factor.  In an
        array of $N$ meta-atoms with fabrication scatter
        $\sigma_{\omega_2}$ on the individual SH-resonance frequencies,
        the collective response is inhomogeneously broadened with an
        effective array linewidth
        $\gamma_2^\text{(array)}\!\simeq\!\sqrt{\gamma_2^{2}+\sigma_{\omega_2}^{2}}$,
        and the cascaded enhancement is set by $Q_2^\text{(array)}$,
        not $Q_2^\text{(element)}$.  For $Q_2^\text{(element)}\sim
        10^{4}$ at $\omega_2/2\pi\!\sim\!400$~THz, preserving the
        cascade requires $\sigma_{\omega_2}/\omega_2\lesssim 10^{-4}$,
        i.e.\ sub-nanometre control of meta-atom dimensions.  Mitigations
        include post-fabrication trimming (electron-beam-induced
        defects, photoresist trimming, thermal annealing), spatially
        averaged readout that recovers only the homogeneous response,
        or operating the SH cavity in an electro-optic feedback loop
        that pulls each element onto resonance.
\end{enumerate}

% =======================================================================
\section{Summary of the principal result}

\begin{equation}
  \boxed{\;\chithree_\text{casc}(\omega_1)
        =-\frac{\eps_0\,\omega_1}{6\,\Dom}\;
          \frac{\bigl|\!\int\!\chitwo_{ijk}(\ri)\,f_{2,i}^*\,f_{1,j}\,f_{1,k}\dthreer\bigr|^{2}}
               {\int|\f_1(\ri)|^{4}\dthreer}
        \;\sim\;-\frac{\omega_1}{6\,n_2^{\,2}\,\Dom}\;|\chitwo|^{2}\;\eta_\text{geom}\,,\;}
\end{equation}
valid for $\gamma_2/2\ll|\Dom|\ll\omega_1$ and $|\beta\,a_1|\ll|\Dom|$,
with $\beta,\alpha_3$ as defined in Section~\ref{sec:tcmt}.  The first
form is exact (within the model); the second is the order-of-magnitude
form one obtains by inserting the energy-normalisation scaling
$|\f_n|^2\sim 1/(\eps_0 n_n^{\,2}V_\text{eff})$ -- only $n_2$ (the SH-mode
effective index) survives, all $n_1$ powers cancelling between
numerator and denominator.

Compared with the original derivation, this version:
\begin{itemize}
  \item uses the full vector mode profile with proper energy
        normalisation (Brillouin) accounting for dispersion, rather
        than a scalar $f(\ri)$ with a single $n$;
  \item carries the tensorial $\chitwo_{ijk}$ (and $\chithree_{ijkl}$)
        and gives an unambiguous overlap integral;
  \item includes the factor $2$ in the DFG polarisation and the factor
        $3$ in SPM, ensuring Manley--Rowe is satisfied and the SHG/DFG
        couplings carry the same coupling constant $\beta$;
  \item carries $\eps_0$ correctly through the polarisation--field
        relation, so that the coupled-mode coefficients
        $\beta\propto\eps_0\omega_1$ and $\alpha_3\propto 3\eps_0\omega_1$
        are dimensionally clean; $\eps_0$ then cancels in the final
        $\chithree_\text{casc}$, as it must for a material parameter;
  \item makes the $1/n_2^{\,2}$ dependence explicit (the $n_1$ powers
        cancel between $|\!\int\!\chitwo f_2^*f_1^2|^2$ and $\int|\f_1|^4$,
        leaving only the SH-mode normalisation), correcting the
        $1/n^3$ scaling of the original note;
  \item includes both the cascaded Kerr \emph{and} the cascaded
        two-photon-loss term, essential for the saturation behaviour;
  \item exposes the bandwidth--enhancement sum rule, the correct
        framing for comparison with bulk (or any non-resonant)
        Kerr-equivalent processes;
  \item adds the symmetry / phase-matching constraints that govern
        $\eta$ in a metasurface platform.
\end{itemize}

Whether the resulting ``fast Kerr from slow SHG'' claim survives depends
on the application: for narrowband, CW-pumped nonlinear photonics
(bistability, frequency-comb interactions, low-threshold parametric
oscillation) the metasurface route can outperform bulk by 1--2 orders of
magnitude; for ultrafast applications the sum rule precludes any
improvement over bulk.

% =======================================================================
\begin{thebibliography}{99}

\bibitem{Ostrovskii1967}
L.~A.~Ostrovskii,
``Self-action of light in crystals,''
\emph{JETP Lett.} \textbf{5}, 272 (1967).

\bibitem{DeSalvo1992}
R.~DeSalvo, D.~J.~Hagan, M.~Sheik-Bahae, G.~Stegeman, E.~W.~Van Stryland,
``Self-focusing and self-defocusing by cascaded second-order effects in KTP,''
\emph{Opt.\ Lett.} \textbf{17}, 28 (1992).

\bibitem{Stegeman1996}
G.~I.~Stegeman, D.~J.~Hagan, L.~Torner,
``$\chi^{(2)}$ cascading phenomena and their applications to all-optical signal processing, mode-locking, pulse compression and solitons,''
\emph{Opt.\ Quantum Electron.} \textbf{28}, 1691 (1996).

\bibitem{Bache2007}
M.~Bache, J.~Moses, F.~W.~Wise,
``Scaling laws for soliton pulse compression by cascaded quadratic nonlinearities,''
\emph{J.~Opt.\ Soc.\ Am.\ B} \textbf{24}, 2752 (2007).

\bibitem{Boyd}
R.~W.~Boyd, \emph{Nonlinear Optics}, 4th ed.\ (Academic Press, 2020).

\bibitem{LandauLifshitz}
L.~D.~Landau and E.~M.~Lifshitz,
\emph{Electrodynamics of Continuous Media}, 2nd ed., \S 80.

\bibitem{SuhWangFan2004}
W.~Suh, Z.~Wang, S.~Fan,
``Temporal coupled-mode theory and the presence of non-orthogonal modes in lossless multimode cavities,''
\emph{IEEE J.~Quantum Electron.} \textbf{40}, 1511 (2004).

\bibitem{Rodriguez2007}
A.~Rodriguez, M.~Solja\v{c}i\'c, J.~D.~Joannopoulos, S.~G.~Johnson,
``$\chi^{(2)}$ and $\chi^{(3)}$ harmonic generation at a critical power in inhomogeneous doubly resonant cavities,''
\emph{Opt.\ Express} \textbf{15}, 7303 (2007).

\bibitem{Liu1999}
X.~Liu, L.~J.~Qian, F.~W.~Wise,
``High-energy pulse compression by use of negative phase shifts produced by the cascaded $\chi^{(2)}{:}\chi^{(2)}$ nonlinearity,''
\emph{Opt.\ Lett.} \textbf{24}, 1777 (1999).

\bibitem{Kristensen2020}
P.~T.~Kristensen, K.~Herrmann, F.~Intravaia, K.~Busch,
``Modeling electromagnetic resonators using quasinormal modes,''
\emph{Adv.\ Opt.\ Photon.} \textbf{12}, 612 (2020).

\bibitem{Sauvan2013}
C.~Sauvan, J.-P.~Hugonin, I.~S.~Maksymov, P.~Lalanne,
``Theory of the spontaneous optical emission of nanosize photonic and plasmon resonators,''
\emph{Phys.\ Rev.\ Lett.} \textbf{110}, 237401 (2013).

\bibitem{Lalanne2018}
P.~Lalanne, W.~Yan, K.~Vynck, C.~Sauvan, J.-P.~Hugonin,
``Light interaction with photonic and plasmonic resonances,''
\emph{Laser \& Photonics Reviews} \textbf{12}, 1700113 (2018).

\bibitem{Hsu2016}
C.~W.~Hsu, B.~Zhen, A.~D.~Stone, J.~D.~Joannopoulos, M.~Solja\v{c}i\'c,
``Bound states in the continuum,''
\emph{Nat.\ Rev.\ Mater.} \textbf{1}, 16048 (2016).

\bibitem{Koshelev2018}
K.~Koshelev, S.~Lepeshov, M.~Liu, A.~Bogdanov, Y.~Kivshar,
``Asymmetric metasurfaces with high-$Q$ resonances governed by bound states in the continuum,''
\emph{Phys.\ Rev.\ Lett.} \textbf{121}, 193903 (2018).

\bibitem{Carletti2018}
L.~Carletti, K.~Koshelev, C.~De~Angelis, Y.~Kivshar,
``Giant nonlinear response at the nanoscale driven by bound states in the continuum,''
\emph{Phys.\ Rev.\ Lett.} \textbf{121}, 033903 (2018).

\bibitem{Koshelev2020}
K.~Koshelev, S.~Kruk, E.~Melik-Gaykazyan, J.-H.~Choi, A.~Bogdanov, H.-G.~Park, Y.~Kivshar,
``Subwavelength dielectric resonators for nonlinear nanophotonics,''
\emph{Science} \textbf{367}, 288 (2020).

\bibitem{Anthur2020}
A.~P.~Anthur, H.~Zhang, R.~Paniagua-Dominguez, D.~A.~Kalashnikov, S.~T.~Ha,
T.~W.~W.~Ma{\ss}, A.~I.~Kuznetsov, L.~Krivitsky,
``Continuous wave second harmonic generation enabled by quasi-bound-states in the continuum on gallium phosphide metasurfaces,''
\emph{Nano Lett.} \textbf{20}, 8745 (2020).

\bibitem{KoshelevKivshar2021}
K.~Koshelev, Y.~Kivshar,
``Dielectric resonant metaphotonics,''
\emph{ACS Photonics} \textbf{8}, 102 (2021).

\bibitem{Bulgakov2017}
E.~N.~Bulgakov, A.~F.~Sadreev,
``Bloch bound states in the radiation continuum in a periodic array of dielectric rods,''
\emph{Phys.\ Rev.\ A} \textbf{96}, 013841 (2017).

\bibitem{DrummondWalls1980}
P.~D.~Drummond, D.~F.~Walls,
``Quantum theory of optical bistability. I.\ Nonlinear polarisability model,''
\emph{J.\ Phys.\ A: Math.\ Gen.} \textbf{13}, 725 (1980).

\bibitem{Keller1996}
U.~Keller, K.~J.~Weingarten, F.~X.~K\"artner, D.~Kopf, B.~Braun, I.~D.~Jung, R.~Fluck, C.~H\"onninger, N.~Matuschek, J.~Aus der Au,
``Semiconductor saturable absorber mirrors (SESAMs) for femtosecond to nanosecond pulse generation in solid-state lasers,''
\emph{IEEE J.~Sel.\ Top.\ Quantum Electron.} \textbf{2}, 435 (1996).

\bibitem{Nikogosyan2005}
D.~N.~Nikogosyan, \emph{Nonlinear Optical Crystals: A Complete Survey} (Springer, 2005).

\end{thebibliography}

\end{document}
